\section{Constraints, Dynamics and Organization}\label{sec:ext}
\IfFileExists{part2/tables/ext_macros.tex}{% Generated by code/verify_part2_ext.py; do not edit. XBOUNDS placeholders: 2.5, 5, 9.
\providecommand{\extXbOne}{2.5}
\providecommand{\extXbTwo}{5}
\providecommand{\extXbThree}{9}
\providecommand{\extXbRetOne}{31}
\providecommand{\extXbRetTwo}{15.5}
\providecommand{\extXbRetThree}{8.6}
}{}% XBOUNDS placeholder values, generated by code/verify_part2_ext.py

The activity model of Section~\ref{sec:tasks} compares the unit costs of in-house work (H), outsourcing (O), automation (A), augmentation (U) and elimination (N) for each activity and process. This section adds what that comparison leaves out. Section~\ref{sec:ext-universe} lets taxes, regulation, ethical and religious commitments, altruism and customer preferences enter the comparison, and defines the set of activities available to automate as four nested layers. Section~\ref{sec:ext-dynamics} lets capability, prices and energy costs move over time, treats the limit of unbounded cognitive capability, and derives the optimal adoption date when costs fall and transitions require parallel running. Section~\ref{sec:ext-competition} uses the entry results of Part I to separate automation that creates value from automation that is defensive and from automation that is feasible but not worth doing. Section~\ref{sec:ext-org} turns the mode assignment into a team size, a hierarchy, a residual organization and a transformation team. Section~\ref{sec:ext-program} sequences a program of transitions under dual-running costs and limited change capacity, Section~\ref{sec:ext-bandwidth} adds risk, learning, failures that processes share through a common model provider, and the organization's limited capacity to make decisions and to review deployments, Section~\ref{sec:ext-shocks} adds leadership changes, vision overhauls, reorganizations and the loss of context, and Section~\ref{sec:ext-breaks} collects the breakpoints. Proofs are in Appendix~\ref{app:proofs-ext}, and \texttt{code/verify\_part2\_ext.py} checks every closed form, threshold and example. Symbols reserved in Part I keep their meanings; ethical and religious commitments and altruism are parameters that the firm or its stakeholders set, and the paper takes no position on them.

\subsection{Setting}\label{sec:ext-setup}

\begin{assumption}[Mode costs]\label{ass:ext-cost}
Activity $j$ has the attributes of Section~\ref{sec:tasks}: annual volume $v_j$, human hours per unit $t_j$, gross wage $w^0_j$ and loaded wage $w_j=w^0_j(1+\tau_L)$, loss per undetected error $L_j$, task length $d_j$ (human hours per instance), review time per reviewed unit $t^R_j$ and augmented time $t^U_j$. For each mode $x\in\mathcal M:=\{\mathrm H,\mathrm O,\mathrm A,\mathrm U,\mathrm N\}$ the unit cost splits into four parts:
\begin{equation}\label{eq:ext-split}
c_{x,j}=\underbrace{w_jt_{x,j}+w^L_jt^L_{x,j}}_{\text{own labor}}+c^{\rm ext}_{x,j}+c^{\rm mach}_{x,j}+c^{K}_{x,j}.
\end{equation}
Own hours per unit are $t_{\mathrm H}=t_j$, $t_{\mathrm O}=\kappa_Ot_j$, $t_{\mathrm U}=t^U_j$, the augmented time of Section~\ref{sec:tasks}, and $t_{\mathrm N}=0$. In mode A the own-labor and loss parts together are the quality cost of Proposition~\ref{prop:tasks-review}, $w_jt_{\mathrm A}+c^{\rm ext}_{\mathrm A}=c^Q_j$: the own hours $t_{\mathrm A}=\xi^*_jt^R_j$ are the review hours at the optimal share $\xi^*_j$ of \eqref{eq:tasks-xi}, and $c^{\rm ext}_{\mathrm A}=e_jL_j[1-\xi^*_j(1+\ln(1/\xi^*_j))]$ is the expected loss from undetected errors when the highest-stakes units are reviewed. Under the failure taxonomy of Section~\ref{sec:tasks-failure}, $c^Q_j$ is the quality cost of Proposition~\ref{prop:tasks-halluc}, with its review-proof floor, and mode A also carries the approval hours $t_j(\mathsf a^*_j)$ and the expected misbehavior loss $c^m_j(\mathsf a^*_j)$ at the optimal autonomy level of Proposition~\ref{prop:tasks-autonomy}; they enter $t_{\mathrm A}$ and $c^{\rm ext}_{\mathrm A}$ in the same way. Licensed sign-off hours $t^L_{\mathrm A}=\mathrm{lic}_j\,t^L_j$ are paid the licensed wage $w^L_j$, with $\mathrm{lic}_j\in\{0,1\}$ and $t^L_x=0$ for $x\neq\mathrm A$. The other purchases and losses are $c^{\rm ext}_{\mathrm O}=p_O+\zeta_O$ and $c^{\rm ext}_{\mathrm N}=L^N_j$, the value lost by dropping the activity. Variable AI spend is $c^{\rm mach}_{\mathrm A}=\pi n$ and $c^{\rm mach}_{\mathrm U}=\pi n_U$. For the AI modes the capital charge is $c^K_{x,j}=(r+\delta_K)K_{x,j}+F_{x,j}\,a(r,T_x)/v_j$, where $K$ is capital per unit of annual volume, $F$ the build and integration cost and $a(r,T)=r/(1-e^{-rT})$ the capital recovery factor; capital used by the non-AI modes is included in their $c^{\rm ext}$. All other parts are zero. In-house work is always feasible and permitted.
\end{assumption}

This reproduces the unit costs of Assumption~\ref{ass:tasks-modes}, with the quality cost $c^Q_j$ split into its labor part, on which payroll taxes and altruism act, and its loss part. The build cost is annualized with the capital recovery factor, which Section~\ref{sec:tasks} allows and which tends to the straight-line charge $1/T_A$ as $r\to0$. The split separates the parts on which taxes, preferences and AI prices act. We call $\mathcal M^{\rm AI}:=\{\mathrm A,\mathrm U,\mathrm N\}$ the AI modes and $\mathcal M^0:=\{\mathrm H,\mathrm O\}$ the non-AI modes.

\begin{assumption}[Policy and preference parameters]\label{ass:ext-theta}
The parameter vector $\Theta$ has five groups.
\begin{enumerate}
\item[(i)] \emph{Taxes.} The payroll loading $\tau_L$ covers payroll taxes and benefits. Profits are taxed at rate $\tau\in[0,1)$ with full loss offset. The share $\psi\in[0,1]$ is the present value of tax deductions per dollar of AI capital (Part I's $\psi$, applied to AI capital); with straight-line deductions over $T_R$ years, $\psi=(1-e^{-rT_R})/(rT_R)$. Under current U.S.\ law, software development is research or experimental expenditure: its domestic part is deducted when paid or incurred, so $\psi=1$, and the part attributable to foreign research is amortized over 15 years from the midpoint of the taxable year \citep{usc2026sec174a,usc2026sec174}. For a build with foreign share $\varsigma_{\rm for}$, $\psi=1-\varsigma_{\rm for}(1-\psi_{15})$, where $\psi_{15}=0.532$ at $r=0.10$ under the statutory midpoint convention (0.518 with continuous straight-line deductions; Proposition~\ref{prop:acct-tax}). Qualified property acquired after 19 January 2025, which includes domestic servers, has a 100\% first-year allowance, so $\psi=1$ for it too, unless the firm elects out; property used predominantly outside the United States is depreciated under the alternative depreciation system and has $\psi<1$ \citep{irs2026notice202611,usc2026sec168}. A hypothetical, deductible automation tax $\tau_A\ge0$ is levied on AI spend and AI capital.
\item[(ii)] \emph{Regulation.} The firm maps each activity's use case to a regulatory class $c_j\in\{0,1,2,3\}$: no obligations, transparency, high risk (for example the areas listed in Annex III of the EU AI Act) and prohibited. The oversight floor of Assumption~\ref{ass:tasks-review} is set by class, $\xi^{\min}_j=\xi^{\min}(c_j)$, with $\xi^{\min}(\cdot)$ nondecreasing, and mode A reviews the optimal share $\xi^*_j$ of Proposition~\ref{prop:tasks-review}, which is at least the floor. Class 3 rules out modes A and U. Under the AI Act as amended in 2026, the high-risk obligations apply from 2 December 2027 for the Annex III areas and from 2 August 2028 for systems under Annex I \citep{eu2026regulation}; a plan whose horizon runs past those dates should treat such use cases as class 2. If $\mathrm{lic}_j=1$, mode A requires sign-off hours and mode N is not permitted. Limits that data protection places on readiness $\iota_j$ act through Section~\ref{sec:tasks}.
\item[(iii)] \emph{Ethics and religion.} A hard exclusion set $\mathcal E\subseteq\mathcal J\times\mathcal M^{\rm AI}$ lists activity-mode pairs ruled out. A soft penalty $\omega_{x,j}\ge0$, in after-tax dollars per unit, applies to AI modes; $\omega_{\mathrm H,j}=\omega_{\mathrm O,j}=0$.
\item[(iv)] \emph{Altruism.} The owner values each own employee hour retained at $\vartheta\ge0$ after-tax dollars.
\item[(v)] \emph{Customer preference.} For a customer-facing activity, AI mode $x$ lowers the firm's zero-spend growth (retention) by $\delta^G_{x,j}$ and its demand pull by the proportion $\delta^\lambda_{x,j}$. Together they are the demand penalty of customer preference for human service.
\end{enumerate}
\end{assumption}

The EU AI Act sets no minimum share of work that a human must review: Article 14 requires that high-risk systems can be overseen, overridden and stopped by natural persons, and Annex III lists use-case areas rather than graded classes \citep{eu2024regulation}. The floor $\xi^{\min}$ and the classes $c_j$ are therefore modeling devices that turn these obligations into cost; the firm sets their values.

\begin{definition}[Effective unit cost and automation advantage]\label{def:ext-effcost}
The owner's objective per unit in mode $x$ is the after-tax cost plus the non-pecuniary terms. Divided by $1-\tau$, it is the effective unit cost
\begin{equation}\label{eq:ext-effcost}
\tilde c_{x,j}(\Theta)=\hat w_jt_{x,j}+\hat w^L_jt^L_{x,j}+c^{\rm ext}_{x,j}+(1+\tau_A)\bigl(c^{\rm mach}_{x,j}+\Upsilon\,c^K_{x,j}\bigr)+\frac{\omega_{x,j}}{1-\tau}+\mathcal C_{x,j},
\end{equation}
where
\begin{equation}\label{eq:ext-wedges}
\hat w_j:=w^0_j(1+\tau_L)-\frac{\vartheta}{1-\tau},\qquad \Upsilon:=\frac{1-\psi\tau}{1-\tau}\ge1,\qquad \mathcal C_{x,j}:=\hat X\bigl(\delta^G_{x,j}+\Delta^*\delta^\lambda_{x,j}\bigr)\frac{S}{v_j},\qquad \hat X:=\frac{X-(1-\psi\tau)k}{1-\tau},
\end{equation}
with $\hat w^L_j$ defined like $\hat w_j$ and $\tau_A$ read as zero for the non-AI modes. In mode A the review share, and under the failure taxonomy the autonomy level, are chosen at the effective wage, so $\hat w_jt_{\mathrm A}+c^{\rm ext}_{\mathrm A}=:\hat c^Q_j$ is the quality cost of Proposition~\ref{prop:tasks-review} (with, under the taxonomy, the approval and misbehavior costs) evaluated at the wage $\hat w_j$; by the envelope theorem, $\partial\hat c^Q_j/\partial\hat w_j=t_{\mathrm A}$, which is $\xi^*_jt^R_j$ in the base case. We assume $\vartheta\le(1-\tau)w^0_j$ and $X\ge(1-\psi\tau)k$, so that $\hat w_j\ge0$ and $\mathcal C\ge0$. Here $X$ is the firm's own growth weight, $X=v'(s)$ by Proposition~\ref{prop:dynamic}, $\Delta^*$ its excess growth, $S$ its revenue and $k$ its capital intensity. Let $\mathcal M_j(t,\Theta)$ be the modes that are feasible and permitted for $j$ at date $t$. The automation advantage per unit is
\begin{equation}\label{eq:ext-B}
\mathcal B_j(t,\Theta):=\min_{x\in\mathcal M^0\cap\mathcal M_j}\tilde c_{x,j}-\min_{x\in\mathcal M^{\rm AI}\cap\mathcal M_j}\tilde c_{x,j},
\end{equation}
which is $-\infty$ if no AI mode is in $\mathcal M_j$. We write $x^0_j$ and $x^1_j$ for the two minimizers.
\end{definition}

Three wedges carry the tax and preference parameters. Altruism lowers the effective cost of an own hour from $w_j$ to $\hat w_j$. The factor $\Upsilon$ is the tax penalty on AI capital, which disappears under full expensing ($\psi=1$). Under current law this is the case for domestic software development and for property with the full first-year allowance. For development performed abroad, $\psi_{15}=0.532$ at $r=0.10$, so at $\tau=0.21$ AI capital costs 12.4\% more than equally productive operating spend. Under the five-year amortization that applied to domestic development in 2022--2024, $\psi=0.787$ and the penalty was 5.7\%. The demand penalty $\mathcal C$ converts a loss of retention or demand pull into dollars per unit at the firm's own growth weight; Proposition~\ref{prop:ext-customer} justifies it.

\subsection{The universe available to automate}\label{sec:ext-universe}

\begin{definition}[Four nested layers]\label{def:ext-universe}
Let $\mathcal M^F_j(t)$ be the modes that are technically feasible for activity $j$ at date $t$: capability (Section~\ref{sec:ext-dynamics}) and readiness at or above its feasibility floor, $\iota_j\ge\underline\iota_j$ (Proposition~\ref{prop:tasks-augment}). Let $\mathcal B^0_j$ be the automation advantage when the variable AI spend and the capital charge of the AI modes are set to zero. Define
\begin{align*}
\mathcal J^{\rm fea}(t)&:=\{j:\mathcal M^F_j(t)\cap\mathcal M^{\rm AI}\neq\emptyset\}&&\text{(feasible)},\\
\mathcal J^{\rm per}(t,\Theta)&:=\{j:\mathcal M_j(t,\Theta)\cap\mathcal M^{\rm AI}\neq\emptyset\}&&\text{(permissible)},\\
\mathcal J^{\rm pre}(t,\Theta)&:=\{j\in\mathcal J^{\rm per}:\mathcal B^0_j(t,\Theta)>0\}&&\text{(preferred)},\\
\mathbb A(t,\Theta):=\mathcal J^{\rm eco}(t,\Theta)&:=\{j\in\mathcal J^{\rm per}:\mathcal B_j(t,\Theta)>0\}&&\text{(economical, the available universe).}
\end{align*}
With the labor-equivalent spend $E_j=v_jw_jt_j$ of Section~\ref{sec:tasks} as fixed weights, the spend share of layer $\bullet$ is $\Sigma^\bullet:=\sum_{j\in\mathcal J^\bullet}E_j/\sum_jE_j$, and the annual value of the universe is $\mathcal V^{\rm eco}:=\sum_jv_j\,(\mathcal B_j)_+$.
\end{definition}

The preferred layer asks whether an AI mode would win if inference and integration were free. It therefore keeps the costs that do not fall with AI prices: the quality cost (review and expected errors), sign-off, soft penalties, altruism and customer preference. $\Sigma^{\rm eco}$ extends the automation curve $S_A$ of Proposition~\ref{prop:tasks-curve} to all AI modes and to the parameters $\Theta$.

\begin{lemma}[Nesting and the waterfall]\label{lem:ext-nested}
\begin{enumerate}
\item[(a)] $\mathbb A\subseteq\mathcal J^{\rm pre}\subseteq\mathcal J^{\rm per}\subseteq\mathcal J^{\rm fea}$. Hence $\Sigma^{\rm eco}\le\Sigma^{\rm pre}\le\Sigma^{\rm per}\le\Sigma^{\rm fea}\le1$, and the spend outside the universe splits into capability-limited, rule-limited, preference-limited and cost-limited parts:
\begin{equation}\label{eq:ext-waterfall}
1-\Sigma^{\rm eco}=(1-\Sigma^{\rm fea})+(\Sigma^{\rm fea}-\Sigma^{\rm per})+(\Sigma^{\rm per}-\Sigma^{\rm pre})+(\Sigma^{\rm pre}-\Sigma^{\rm eco}).
\end{equation}
\item[(b)] $\mathcal J^{\rm pre}$ does not depend on the price of inference $\pi$, on build costs $F$, on capital per unit $K$, or on $\tau_A$. As these AI costs fall to zero, $\mathbb A$ grows to $\mathcal J^{\rm pre}$. No fall in AI prices or energy costs can therefore raise $\Sigma^{\rm eco}$ above $\Sigma^{\rm pre}$.
\end{enumerate}
\end{lemma}

The four gaps behave differently over time. The capability gap closes as the horizon grows (Proposition~\ref{prop:ext-feasible}) and the cost gap as prices fall. The rule and preference gaps close only if $\Theta$ changes. A firm that tracks only the economical layer cannot tell a universe that will grow on its own from one that is held back by rules or preferences.

\begin{proposition}[Comparative statics of the available universe]\label{prop:ext-statics}
\begin{enumerate}
\item[(a)] (Unconditional.) $\mathbb A(t,\Theta)$, $\Sigma^{\rm eco}$ and $\mathcal V^{\rm eco}$ are nonincreasing in the automation tax $\tau_A$, the inference price $\pi$, the build cost $F_A$, the capital per unit $K_A$, each oversight floor $\xi^{\min}(c)$, each penalty $\omega_{x,j}$, the demand terms $\delta^G$ and $\delta^\lambda$, the exclusion set $\mathcal E$ (in set inclusion), a higher class $c_j$ and licensing. They are nondecreasing in the deductibility $\psi$ and in readiness $\iota$ (through the raw error rate $e_j(\iota)$, and hence $\xi^*_j$ and $c^Q_j$).
\item[(b)] (Conditional.) Where the minimizers $x^0_j$ and $x^1_j$ are unique,
\begin{align}
\frac{\partial\mathcal B_j}{\partial\tau_L}&=w^0_j\bigl(t_{x^0_j}-t_{x^1_j}\bigr)+w^{L0}_j\bigl(t^L_{x^0_j}-t^L_{x^1_j}\bigr),\qquad
\frac{\partial\mathcal B_j}{\partial\vartheta}=-\frac{\bar t_{x^0_j}-\bar t_{x^1_j}}{1-\tau},\label{eq:ext-dBdL}\\
\frac{\partial\mathcal B_j}{\partial\tau}&=-\frac{(1+\tau_A)(1-\psi)\,c^K_{x^1_j}+\omega_{x^1_j}+\vartheta\bigl(\bar t_{x^0_j}-\bar t_{x^1_j}\bigr)+\delta_{x^1_j}\frac{S}{v_j}\bigl(X-(1-\psi)k\bigr)}{(1-\tau)^2},\label{eq:ext-dBdtau}
\end{align}
where $\bar t_x:=t_x+t^L_x$ is total own hours, with $t_{\mathrm A}$ the own hours of mode A at its optimal review share (and autonomy level), by the envelope theorem, $\delta_x:=\delta^G_x+\Delta^*\delta^\lambda_x$, and $X$ and $k$ are held fixed. A payroll tax enlarges the universe, and altruism shrinks it, exactly when the AI mode uses fewer own hours than the best non-AI mode. When that mode is outsourcing with fewer own hours than the AI mode, both signs reverse. The profit tax shrinks the universe when the AI mode displaces own hours, $\omega\ge0$ and $X\ge(1-\psi)k$. It is neutral when $\psi=1$ and there are no non-pecuniary terms.
\item[(c)] (Equivalences.) The payroll loading $\tau_L$ and altruism $\vartheta$ enter $\mathcal B_j$ only through $\hat w_j$. Altruism $\vartheta$ is equivalent to a payroll subsidy of $\vartheta/(1-\tau)$ per hour. A deductible tax $T_D$ on each own hour displaced, $T_D(t_j-\bar t_{x,j})$ in mode $x$, gives the same $\mathcal B_j$ as altruism $\vartheta=(1-\tau)T_D$. Mode choices therefore cannot distinguish displacement taxes, altruism and payroll subsidies; their effects on revenue and welfare differ.
\end{enumerate}
\end{proposition}

\noindent\emph{Intuition.} A parameter that raises the cost of every AI mode at least as much as the cost of every non-AI mode can only shrink the universe. Most parameters are of this kind (part a). Payroll taxes and altruism are not. They fall on hours, so their effect depends on the hours that the best AI mode displaces relative to the best non-AI mode, and that alternative is often outsourcing rather than in-house work. A payroll tax then favors the vendor, whose staff it does not tax. Altruism toward employees penalizes outsourcing most, because outsourcing removes every own job. Because the monotone results need no differentiability, they also hold for the discrete changes the engine evaluates; they are monotone comparative statics in the sense of \citet{milgrom1994monotone}. \citet{acemoglu2020does} argue that the U.S.\ tax code favors automation by taxing labor more heavily than capital. Part (b) adds that, at the level of the firm's choice, this holds when the alternative is in-house work and can fail when it is outsourcing.

\begin{table}[t]
\centering
\caption{How each parameter moves the effective unit costs and the available universe}\label{tab:ext-statics}
\footnotesize
\setlength{\tabcolsep}{4pt}
\begin{tabular}{@{}p{3.3cm}p{4.3cm}p{1.5cm}p{5.6cm}@{}}
\toprule
Parameter & Enters $\tilde c$ through & Effect on $\Sigma^{\rm eco}$ & Condition\\
\midrule
Payroll loading $\tau_L$ & $\hat w$ on own hours & $+$ / $-$ & sign of hours displaced vs.\ best non-AI mode\\
Profit tax $\tau$, deductibility $\psi$ & $\Upsilon$ on AI capital; $1/(1-\tau)$ on preferences & $-$ ($\tau$), $+$ ($\psi$) & $\tau$: AI mode displaces hours; neutral if $\psi=1$ and no preferences\\
Automation tax $\tau_A$ & AI spend and capital & $-$ & none\\
Oversight floor $\xi^{\min}(c)$ & quality cost of mode A & $-$ & binds where it exceeds the economic and reliability shares\\
Licensing & sign-off hours; N not permitted & $-$ & none\\
Exclusions $\mathcal E$, class 3 & removes modes & $-$ & set inclusion\\
Ethics penalty $\omega$ & $\omega/(1-\tau)$ on AI modes & $-$ & none\\
Altruism $\vartheta$ & $-\vartheta/(1-\tau)$ per own hour & $-$ / $+$ & opposite to $\tau_L$\\
Customer preference $\delta^G,\delta^\lambda$ & $\mathcal C=\hat X\delta S/v$ & $-$ & weighted by the marginal $q$ of revenue\\
Readiness $\iota$ & $e_j(\iota)$, hence $\xi^*_j$ and $c^Q_j$; feasibility & $+$ & none\\
\bottomrule
\end{tabular}
\end{table}

\begin{proposition}[The spend share is fragile, the value is robust]\label{prop:ext-margin}
Let a scalar component $\theta$ of $\Theta$, or an AI price, move every $\mathcal B_j$ continuously, with derivative $b_j:=\partial\mathcal B_j/\partial\theta$ away from crossings. The advantage is affine near $\theta_0$, $\mathcal B_j(\theta)=\mathcal B_j(\theta_0)+b_j(\theta-\theta_0)$, in $\tau_A$, $\omega$, $\delta^G$, $\delta^\lambda$, $\pi$, $F_A$ and $K_A$ while the minimizers do not change, and also in $\tau_L$ and $\vartheta$ where the review share of mode A is set by the reliability requirement or the floor rather than by the wage.
\begin{enumerate}
\item[(a)] For a finite set of activities, $\Sigma^{\rm eco}$ is a step function of $\theta$: it jumps by $E_j/\sum_iE_i$ when $\mathcal B_j$ crosses zero. $\mathcal V^{\rm eco}$ is continuous, piecewise affine where the $\mathcal B_j$ are affine, and away from crossings
\begin{equation}\label{eq:ext-envelope}
\frac{d\mathcal V^{\rm eco}}{d\theta}=\sum_{j:\mathcal B_j>0}v_j\,b_j .
\end{equation}
Activities at the margin contribute nothing to this derivative.
\item[(b)] For a continuum of activities, with each $\mathcal B_j$ affine in $\theta$ and $\mathcal B_j(\theta_0)$ having, under the spend weights, a density conditional on $b_j$ that is continuous at zero,
\begin{equation}\label{eq:ext-density}
\frac{d\Sigma^{\rm eco}}{d\theta}=f_{\mathcal B}(0)\;\mathbb E^E\bigl[b\mid\mathcal B=0\bigr],
\end{equation}
where $f_{\mathcal B}$ is the spend-weighted density of the automation advantage.
\item[(c)] If $\theta$ enters only the costs of AI modes, and does so affinely (as $\tau_A$, $\pi$, $F_A$, $K_A$, $\omega$, $\delta^G$ and $\delta^\lambda$ do), then $\mathcal V^{\rm eco}$ is convex in $\theta$. A first-order estimate therefore understates the value gained when AI becomes cheaper and overstates the value lost when it becomes dearer.
\end{enumerate}
\end{proposition}

\noindent\emph{Intuition.} The spend share responds through the activities at the margin, where the density of the advantage can be high. The value responds only through activities already inside, because a marginal activity is worth almost nothing to automate. A tornado on the spend share can therefore swing widely while the value hardly moves, and the engine should rank sensitivities by value.

\begin{proposition}[Customer preference and the marginal $q$ of revenue]\label{prop:ext-customer}
Let automating customer-facing activity $j$ in mode $x$ raise the pre-tax margin intercept by $dm_0$, lower zero-spend growth by $\delta^G$ and lower demand pull by the proportion $\delta^\lambda$, all small, with $\bar\sigma=\infty$.
\begin{enumerate}
\item[(a)] Under Remark~\ref{rem:taxes} and Corollary~\ref{cor:priorities}, the after-tax maximal score changes to first order by $(1-\tau)\,dm_0-\bigl[X-(1-\psi\tau)k\bigr]\bigl(\delta^G+\Delta^*\delta^\lambda\bigr)$. With $X=v'(s)$ this is the firm's own per-period criterion (Proposition~\ref{prop:dynamic}(c)). It is positive if and only if the pecuniary saving per unit exceeds $\mathcal C_{x,j}$ of \eqref{eq:ext-wedges}.
\item[(b)] In the small-share limit of Proposition~\ref{prop:dynamic}(d), with zero-spend growth $G_i$ in place of $G$ and $\tau=0$,
\begin{equation}\label{eq:ext-dv0}
dv_0=\frac{dm_0-(v_0-k)\bigl(\delta^G+\Delta^\circ\delta^\lambda\bigr)}{r-g^\circ},
\end{equation}
so the criterion of (a) with $X=v_0$ is exact to first order in firm value.
\end{enumerate}
\end{proposition}

\noindent\emph{Intuition.} A saving counts once, but a loss of retention costs the firm its marginal value of revenue, the decision weight $X$, which the data do not pin down: Part I's spend-based estimate of 7 to 13 is withdrawn, and bounds by method run from below one to about twelve for $X-k$. Customer-facing automation must therefore clear a bar that rises with $X$, and that is higher than back-office automation's with the same saving whenever $X$ exceeds $(1-\psi\tau)k$. Customers' resistance to algorithms is documented \citep{dietvorst2015algorithm,longoni2019resistance}, and disclosure that a sales call came from a chatbot reduced purchases sharply in a field experiment \citep{luo2019machines}. The size of $\delta^G$ for a given activity is an assumption that the engine should vary.

\begin{corollary}[An activity-level accountability floor]\label{cor:ext-floor}
Let capability be unbounded (Proposition~\ref{prop:ext-agi}), so that every permissible activity can run in mode A with its review fraction at the floor, and let impermissible activities stay in their best non-AI mode. Under full automation the residual human cost per dollar of revenue is
\begin{equation}\label{eq:ext-amin}
\frac1S\sum_{j\in\mathcal J^{\rm per}}v_j\Bigl[w_j\xi^{\min}(c_j)\,t^R_j+\mathrm{lic}_j\,w^L_jt^L_j\Bigr]+\frac1S\sum_{j\notin\mathcal J^{\rm per}}v_j\,c_{x^0_j,j},
\end{equation}
which splits into a floor $c_h^{\min}(\Theta)$ on human service cost (activities in cost of revenue) and a floor on operating cost that is part of $a_{\min}$. Corollary~\ref{cor:ceiling} then holds with $c_h=c_h^{\min}(\Theta)$, and $\partial R_{\max}/\partial\xi^{\min}(c)=-(1-\tau)\sum_{j:c_j=c}v_jw_jt^R_j/S$ over the activities whose floor binds.
\end{corollary}

The accountability floor of Part I thus has an activity-level counterpart that the engine can compute from oversight floors, licensing and exclusions. Under the failure taxonomy of Section~\ref{sec:tasks-failure}, the bracket in \eqref{eq:ext-amin} also includes the wage cost of the human hours $t_j(\mathsf a^*_j)$ at the optimal autonomy level, and any loss floor $c_{\rm err}$ of \eqref{eq:tasks-cerr} that remains enters the ceiling as a separate cost.

\begin{definition}[Tornado and Shapley decomposition]\label{def:ext-tornado}
For each parameter $\theta_k$ with an engine range $[\theta_k^{\rm lo},\theta_k^{\rm hi}]$, the tornado entries are $\Sigma^{\rm eco}$ and $\mathcal V^{\rm eco}$ at $(\theta_k^{\rm lo},\Theta^0_{-k})$ and $(\theta_k^{\rm hi},\Theta^0_{-k})$, sorted by the range of $\mathcal V^{\rm eco}$. For a scenario that moves the parameter groups $1,\dots,n_g$ from $\Theta^0$ to $\Theta^1$, let $f(T)$ be $\Sigma^{\rm eco}$ (or $\mathcal V^{\rm eco}$) with the groups in $T$ at $\Theta^1$ and the others at $\Theta^0$. The Shapley attribution of group $k$ is
\begin{equation}\label{eq:ext-shapley}
\mathrm{Sh}_k:=\sum_{T\subseteq\{1..n_g\}\setminus\{k\}}\frac{|T|!\,(n_g-|T|-1)!}{n_g!}\bigl[f(T\cup\{k\})-f(T)\bigr].
\end{equation}
\end{definition}

\begin{proposition}[Sensitivity decomposition]\label{prop:ext-shapley}
\begin{enumerate}
\item[(a)] $\sum_k\mathrm{Sh}_k=f(\{1..n_g\})-f(\emptyset)$, so the attributions add up to the total change exactly.
\item[(b)] If $f$ is additively separable across groups, $\mathrm{Sh}_k$ equals the one-at-a-time change $f(\{k\})-f(\emptyset)$.
\item[(c)] Otherwise the one-at-a-time changes need not add up to the total. Because $\mathcal B_j$ is a difference of minima, constraints interact (for example, an oversight floor matters only where readiness has made mode A the best AI mode), so interactions are the rule.
\end{enumerate}
\end{proposition}

The decomposition needs $2^{n_g}$ evaluations, 64 for the five groups of Assumption~\ref{ass:ext-theta} plus AI prices, which the engine can afford for every scenario comparison, for example an EU against a U.S.\ parameter set.

\paragraph{Example.} Take customer-support tickets: $v=10^6$ a year, $t=0.25$ hours, a gross wage of \$40 an hour and the private-industry loading $\tau_L=0.43$, the ratio of benefits to wages in the June 2026 Employer Costs for Employee Compensation (\$14.07 against \$32.82 an hour). Then $c_{\mathrm H}=\$14.30$. Outsourcing at $p_O=\$7.00$, with $\kappa_O=0.10$ and $\zeta_O=\$0.60$, costs \$9.03. Automation with inference \$0.40, a \$1.5 million build over three years ($r=0.10$), review time $t^R=t$, a raw error rate $e=0.10$, a reliability target $\bar q=0.915$ and a loss of \$17.69 per undetected error costs \$4.12 at $\tau=0.21$ and $\psi=1$, a domestic build deducted when incurred. The reliability target sets the review share at $\xi^*=1-\bar\varepsilon/e=0.15$ (the economic share is 0.0003), so $c^Q=0.15\times\$14.30+\$1.00=\$3.15$. Augmentation that raises issues resolved per hour by 15\%, the average gain that \citet{brynjolfsson2025generative} report for support agents ($t^U=t/1.15$), costs \$12.53. The best non-AI mode is outsourcing, and $\mathcal B=\$4.91$ a ticket, or \$4.91 million a year. The activity leaves the universe at an oversight floor of $\xi^{\min*}=0.55$. Because forced review also catches errors, this lies above the bound $\xi^*+\mathcal B/(\hat wt^R)=0.49$ of Proposition~\ref{prop:ext-breaks}(d), which would hold with equality if the error cost did not fall with review. Because the comparator is outsourcing, $\partial\mathcal B/\partial\tau_L=40\,(0.025-0.0375)=-0.50$, and altruism of \$10 an hour \emph{raises} $\mathcal B$ to \$5.06. If data protection rules out the vendor, the comparator is in-house work, $\partial\mathcal B/\partial\tau_L=+8.5$, and the same altruism lowers $\mathcal B$ from \$10.18 to \$7.49. For a firm with \$500 million of revenue and $k=0$, the automation pays only if it costs less than $\delta^{G*}=\mathcal B v/(\hat XS)$ of retention, which is inversely proportional to the growth weight: \extXbRetOne, \extXbRetTwo{} and \extXbRetThree{} basis points at illustrative weights of \extXbOne, \extXbTwo{} and \extXbThree. The sign of the decision does not flip across these weights; the tolerated loss rises without bound as $X$ falls toward zero, so below \extXbOne{} it exceeds \extXbRetOne{} basis points. % XBOUNDS: placeholder weights (macros extXb*, set in code/verify_part2_ext.py)

Table~\ref{tab:ext-tornado} shows the tornado layout for a synthetic population of 500 activities. The population is drawn by the verification script from stated ranges; it is not data. Its baseline deductibility is current law for domestic development, $\psi=1$, and the $\psi$ row moves the foreign share of development from 1 ($\psi_{15}=0.532$) to 0. Mode A's quality cost follows Proposition~\ref{prop:tasks-review}, with review times, raw error rates, losses per error and reliability targets drawn for each activity. The baseline waterfall is $\Sigma^{\rm fea}=0.862$, $\Sigma^{\rm per}=0.824$, $\Sigma^{\rm pre}=0.776$ and $\Sigma^{\rm eco}=0.742$. In this population outsourcing is available for 55\% of activities, and augmentation is often the best AI mode. A higher payroll loading therefore \emph{lowers} the spend share, by moving activities from augmentation to outsourcing, while it \emph{raises} the value, through the activities whose comparator is in-house work. Altruism does the reverse. For a scenario that moves all six groups of the Shapley decomposition at once, $\Sigma^{\rm eco}$ rises by 0.8 points. The six one-at-a-time changes sum to 5.1 points, while the Shapley attributions sum to exactly 0.8. Under the 2022--2024 rule ($\psi=0.787$), the baseline spend share was 0.735, the scenario raised it by 1.5 points and the one-at-a-time changes summed to 6.9; the ranking of drivers in Table~\ref{tab:ext-tornado} is unchanged.

\begin{table}[t]
\centering
\caption{Tornado for a synthetic population of 500 activities (spend share $\Sigma^{\rm eco}$ and value $\mathcal V^{\rm eco}$ as \% of baseline spend; baseline 0.742 and 25.0\%)}\label{tab:ext-tornado}
\footnotesize
\begin{tabular}{@{}lccccc@{}}
\toprule
Parameter & Range (low, high) & $\Sigma^{\rm eco}$ low & $\Sigma^{\rm eco}$ high & $\mathcal V^{\rm eco}$ low & $\mathcal V^{\rm eco}$ high\\
\midrule
Inference price (multiple) & 0.1, 3.0 & 0.752 & 0.693 & 27.5 & 20.2\\
Altruism $\vartheta$ (\$ per hour) & 0, 15 & 0.742 & 0.777 & 25.0 & 19.7\\
Payroll loading $\tau_L$ & 0.25, 0.55 & 0.767 & 0.713 & 22.5 & 25.9\\
Readiness: raw error rate $e_j$ (multiple) & 0.6, 1.4 & 0.748 & 0.735 & 27.3 & 24.0\\
Customer penalty (multiple) & 0, 3 & 0.754 & 0.710 & 26.3 & 23.4\\
Build cost (multiple) & 0.5, 2.0 & 0.746 & 0.729 & 25.8 & 23.9\\
Automation tax $\tau_A$ & 0, 0.5 & 0.742 & 0.726 & 25.0 & 23.2\\
Ethics penalty (scale, sensitive activities) & 0, 0.5 & 0.742 & 0.738 & 25.0 & 24.6\\
Oversight floor, high-risk class & 0.05, 0.60 & 0.742 & 0.742 & 25.1 & 24.7\\
Profit tax $\tau$ & 0.15, 0.30 & 0.742 & 0.740 & 25.1 & 24.9\\
Deductibility $\psi$ (foreign share 1, 0) & 0.532, 1.0 & 0.735 & 0.742 & 24.9 & 25.0\\
\bottomrule
\end{tabular}

\smallskip
\begin{minipage}{0.95\linewidth}\footnotesize
Notes: synthetic population and ranges defined in \texttt{code/verify\_part2\_ext.py}; rows sorted by the range of $\mathcal V^{\rm eco}$. The table illustrates the layout and Propositions~\ref{prop:ext-statics} and \ref{prop:ext-margin}; it is not an estimate.
\end{minipage}
\end{table}

\subsection{Capability, prices and adoption timing}\label{sec:ext-dynamics}

\begin{assumption}[Capability path]\label{ass:ext-capability}
\begin{enumerate}
\item[(i)] (Horizon.) The 50\% time horizon, the human length of tasks that AI completes with probability one half, is $H(t)=H_0\,2^{(t-t_0)/T_d}$. \citet{kwa2025measuring} estimate a doubling time of 207 days (95\% interval 166 to 240) over 2019 to 2025, and their task suites predict 143 days from models released in 2024 (a prediction without an interval). We treat $T_d$ as a range from 143 to 240 days. The horizon is measured on software tasks; applying it to other activities is an assumption. In the plateau scenario the horizon stops growing at a date $t_p$, so that $H(t)=H(\min\{t,t_p\})$.
\item[(ii)] (Reliability.) On a task of length $d$, AI succeeds with probability $p(d,t)=\bigl[1+(d/H(t))^{1/\eta_r}\bigr]^{-1}$, logistic in log length. The horizon at reliability $\bar q$ is therefore $H(t)/\chi(\bar q)$, with $\chi(\bar q)=[\bar q/(1-\bar q)]^{\eta_r}$. Horizons at 80\% are 4 to 6 times shorter than at 50\% \citep{kwa2025measuring}, so $\eta_r\in[1.00,1.29]$, with $\eta_r=\ln5/\ln4=1.16$ at the midpoint. Above 80\% the form is an extrapolation.
\item[(iii)] (Context.) Task length splits into a generic part and a firm-specific context part, $d_j=d^g_j+d^c_j$, with context share $s^{\rm ctx}_j=d^c_j/d_j$. An organizational context engine (enterprise data platform, retrieval, long-term memory) with coverage $C_j\in[0,1]$, used at readiness $\iota_j$, supplies the share $\iota_jC_j$ of the context. The task the AI faces then has length $d^g_j+(1-\iota_jC_j)d^c_j$, equivalent to the multiplier $\kappa_{{\rm ctx},j}=1/(1-s^{\rm ctx}_j\iota_jC_j)$ on the horizon.
\item[(iv)] (Physical activities.) An activity flagged physical ($\mathrm{ph}_j=1$) faces the lagged horizon $H(t-L_W)$, where $L_W\ge0$ is the world-model lag (assumed).
\item[(v)] (Prices.) Once an activity is feasible, its AI spend per unit falls at the rate $\delta_\pi$ observed for a fixed level of performance: about 5 to 10 times a year for models on the cost-performance frontier \citep{gundlach2025price}, and 9 to 900 times a year depending on the threshold \citep{epoch2025inference}. The cost of running the frontier model on a benchmark has risen \citep{gundlach2025price}, because the tokens used per task have risen, not because per-token prices have, so an activity that has only just become feasible starts at the frontier cost.
\end{enumerate}
Define, as in Section~\ref{sec:tasks}, the effective capability $\mathcal H:=H(t-L_W\mathrm{ph}_j)$ that activity $j$ faces at date $t$ (written $\mathcal H_j(t)$ where the dependence matters) and its reliability-adjusted difficulty $\bar d_j:=d_j\chi(\bar q^A_j)/\kappa_{{\rm ctx},j}$, where $\bar q^A_j$ is the raw reliability at which mode A is viable. Mode A is capability-feasible for $j$ at $t$ if and only if $\mathcal H\ge\bar d_j$, and mode U under the same test with $d_j$ replaced by the length of the sub-tasks it covers. The raw error rate of Assumption~\ref{ass:tasks-review} is $e_j(t)=1-[1+(d_j/(\kappa_{{\rm ctx},j}\mathcal H_j(t)))^{1/\eta_r}]^{-1}$, and the capability margin $\ln(\mathcal H/\bar d_j)$ is the one that enters the success probability \eqref{eq:tasks-success}.
\end{assumption}

\begin{proposition}[Feasibility over time]\label{prop:ext-feasible}
Let $n^{\rm dbl}_j:=\log_2(\bar d_j/H_0)$ be the doublings that activity $j$ still needs.
\begin{enumerate}
\item[(a)] Activity $j$ becomes capability-feasible at $T^F_j=t_0+T_dn^{\rm dbl}_j+L_W\mathrm{ph}_j$, or now if that date has passed.
\item[(b)] (Reliability lag.) Raising the required reliability from $\bar q$ to $\bar q'$ delays feasibility by $\eta_rT_d\log_2\{[\bar q'/(1-\bar q')]/[\bar q/(1-\bar q)]\}$. From 50\% to 80\%, 95\% and 99\% this is 2.32, 4.93 and 7.70 doublings: 1.32, 2.79 and 4.36 years at $T_d=207$ days, and 3.01 years for 99\% at 143 days.
\item[(c)] (Context.) Raising $\kappa_{\rm ctx}$ from $\kappa$ to $\kappa'$ advances feasibility by $T_d\log_2(\kappa'/\kappa)$. The factor is at most $1/(1-s^{\rm ctx}_j)$. With $s^{\rm ctx}=0.5$, raising $\iota C$ from 0.2 to 0.9 raises $\kappa_{\rm ctx}$ by the factor 1.64: 0.71 doublings, or 0.40 years.
\item[(d)] (The feasible share.) Let $F_g$ be the spend-weighted distribution of $n^{\rm dbl}_j$ within group $g$ (cognitive or physical). Then the capability-feasible spend share is $\sum_gE_g\,F_g\bigl((t-t_0-L_W\mathrm{ph}_g)/T_d\bigr)$. If $n^{\rm dbl}\sim\mathcal N(\bar n^{\rm dbl},\varsigma_{\rm dbl}^2)$ within a group, the share follows a probit curve in calendar time with median date $t_0+T_d\bar n^{\rm dbl}+L_W\mathrm{ph}_g$, and it takes $2z_{0.9}\,\varsigma_{\rm dbl}T_d$ years to go from 10\% to 90\% feasible: 4.36 years for $\varsigma_{\rm dbl}=3$ doublings at 207 days, and 3.01 years at 143 days.
\item[(e)] (Scaling.) With doubling time $T_d'$, the feasible share at $t$ equals the share with $T_d$ at the date $t_0+(T_d/T_d')(t-t_0)$. Multiplying every $\kappa_{\rm ctx}$ by $\kappa'/\kappa$ moves every date earlier by $T_d\log_2(\kappa'/\kappa)$, and a uniform reliability requirement moves them later by the lag in (b).
\item[(f)] (Plateau.) If the horizon stops growing at $t_p$, activity $j$ ever becomes capability-feasible if and only if $\bar d_j\le H(t_p)$, that is, $n^{\rm dbl}_j\le(t_p-t_0)/T_d$, at the date of (a); for a physical activity the plateau binds from $t_p+L_W$. The feasible share of (d) stops rising at $\sum_gE_gF_g((t_p-t_0)/T_d)<1$, so the capability gap of the waterfall \eqref{eq:ext-waterfall} no longer closes with time. An activity above the plateau becomes feasible only if context engines raise $\kappa_{{\rm ctx},j}$ by the factor $\bar d_j/H(t_p)$, or if a lower reliability requirement lowers $\bar d_j$.
\end{enumerate}
\end{proposition}

\noindent\emph{Intuition.} With an exponential horizon, calendar time is a log scale on task length. Every doubling of required length or of the required odds of success costs one doubling time, and context engines buy doublings. The process percolation of Section~\ref{sec:tasks} uses these dates: a process becomes feasible at the latest $T^F_j$ on its critical path. Organizational context is worth roughly months of frontier progress, but it is available now, specific to the firm and therefore a possible scarce input in the sense of Proposition~\ref{prop:entry}(c). Under a plateau it becomes the only lever: if the median activity needs four more doublings from early 2025 and $\varsigma_{\rm dbl}=3$, a plateau in 2027 leaves the capability-feasible share at 39\% for good, and the rest can be reached only through context and readiness.

\begin{proposition}[Inference prices and energy]\label{prop:ext-energy}
Write AI spend per unit as an energy part and a part for hardware, facilities and margins, $c^{\rm mach}=c^{\rm en}+c^{\rm hw}$, with energy share $s_E:=c^{\rm en}/c^{\rm mach}$. Let energy cost per unit fall at the rate $\delta_{\rm en}$ at constant electricity prices, and the rest at $\delta_{\rm hw}$.
\begin{enumerate}
\item[(a)] A proportional fall $\varepsilon_E$ in electricity prices lowers AI spend per unit by the fraction $s_E\varepsilon_E$, which equals $-\ln(1-s_E\varepsilon_E)/\delta_\pi$ years of the price trend. Energy is about \$0.6 billion of the \$8.5 billion annualized cost of a one-gigawatt AI data center, about 7\% \citep{michael2026servers}; its share of a price that includes margins is lower. At $s_E=0.07$, free energy is worth 16.5 days of a fivefold annual price decline, and 11.5 days of a tenfold one.
\item[(b)] The energy share evolves as $\dot s_E=s_E(1-s_E)(\delta_{\rm hw}-\delta_{\rm en})$. It tends to one if hardware and the other costs fall faster than energy costs, and only then does the price of energy become the floor of inference cost.
\item[(c)] Energy prices cannot change $\mathcal J^{\rm fea}$, $\mathcal J^{\rm per}$ or $\mathcal J^{\rm pre}$. Their only effect on the universe is to move activities from the cost-limited gap $\mathcal J^{\rm pre}\setminus\mathbb A$ into $\mathbb A$. We do not model the channel through which energy might limit the compute used to train models, and with it $T_d$.
\item[(d)] At given output prices, the full-automation ceiling of Corollary~\ref{cor:ceiling} rises by $(1-\tau)s_E\varepsilon_Ec_c$: with $\tau=0.18$ and a 90\% fall in electricity prices, 0.31 points at Part I's central compute share at 2025 prices ($c_c=0.06$), and 0.62, 0.41 and 0.26 points at its central shares for 2028, 2031 and 2036 (0.12, 0.08 and 0.05). Under free entry the gain passes to customers (Proposition~\ref{prop:entry}(b)).
\item[(e)] The same formulas apply to physical activities, whose energy share is plausibly larger (assumed).
\end{enumerate}
\end{proposition}

\noindent\emph{Intuition.} Cheap energy, for example from fusion, is a second-order effect for cognitive automation. It lowers a cost that falls much faster for other reasons, and it leaves capability, rules, preferences and accountability untouched. It matters for physical work and in a long run in which hardware has become so cheap that energy is most of what remains.

\begin{proposition}[The limit of unbounded cognitive capability]\label{prop:ext-agi}
Let $H(t)\to\infty$ for cognitive work, with $\Theta$, readiness and the world-model lag fixed.
\begin{enumerate}
\item[(a)] For every cognitive activity the capability test holds and $p(d_j,t)\to1$, so the raw error rate vanishes. $\mathcal J^{\rm fea}$ tends to the activities that meet their readiness floor $\underline\iota_j$ and are cognitive or feasible through world models.
\item[(b)] As the raw error rate vanishes, the review share tends to the floor $\xi^{\min}(c_j)$ and the quality cost to $\xi^{\min}(c_j)\hat w_jt^R_j$ (Proposition~\ref{prop:tasks-review}(d)), so the effective cost of mode A tends to
\begin{equation}\label{eq:ext-cagi}
\tilde c^{\,\infty}_{\mathrm A,j}=\hat w_j\xi^{\min}(c_j)t^R_j+\mathrm{lic}_j\hat w^L_jt^L_j+(1+\tau_A)\bigl(c^{\rm mach,\infty}_{\mathrm A,j}+\Upsilon c^K_{\mathrm A,j}\bigr)+\frac{\omega_{\mathrm A,j}}{1-\tau}+\mathcal C_{\mathrm A,j},
\end{equation}
which does not depend on $H$. This assumes that every failure rate falls to zero with capability; if hallucination or misbehavior rates do not, the review-proof and misbehavior losses of Proposition~\ref{prop:tasks-cerr} remain in the limit.
\item[(c)] The limit universe therefore depends only on $\Theta$, readiness, the cost floors and physical feasibility. It is all of spend if and only if, for every activity, (i) the activity is cognitive or feasible through world models, (ii) $\iota_j\ge\underline\iota_j$, (iii) no exclusion, class-3 status or licensing rules out every AI mode, and (iv) the best limit AI cost, including oversight, sign-off, penalties, altruism and customer preference, is below the best non-AI cost.
\item[(d)] Under full automation of the limit universe, the ceiling of Corollary~\ref{cor:ceiling} holds with the floors of Corollary~\ref{cor:ext-floor}, together with any loss floor $c_{\rm err}$ of \eqref{eq:tasks-cerr} that remains, evaluated at the catch rates of the review that remains (Proposition~\ref{prop:tasks-cerr}(c)), and does not depend on $H$. Once the last cognitive activity is feasible, faster capability growth no longer raises it.
\end{enumerate}
\end{proposition}

Unbounded capability does not remove the accountability floor of Part I; it makes that floor, together with rules, preferences, data readiness, physical work and the cost of compute, the binding constraint.

\begin{proposition}[Optimal adoption date]\label{prop:ext-timing}
Consider a process with volume $v(t)$, best non-AI unit cost $c_0(t)$ (which may fall as vendors adopt AI) and variable AI unit cost $c^{\rm var}_A(t)$ paid as used, so that the savings flow is $W(t):=v(t)[c_0(t)-c^{\rm var}_A(t)]$ a year, the saving $W$ of Assumption~\ref{ass:tasks-transition}. Adopting at date $T$ costs a one-time transition cost $\mathcal T(T)$: the build cost $F_A(T)$, capital $K_Av$, and the present values of the parallel run and of the stranded costs of Section~\ref{sec:tasks}. Savings ramp up as $\varrho(\tau)$ over the parallel-run period $T_{\rm par}$. With the execution risk of Assumption~\ref{ass:tasks-transition}, read $W$ as the risk-adjusted saving $\varpi_p\bar b_pW$ and include in $\mathcal T$ the expected stranded and rollback terms of \eqref{eq:tasks-npv}, where an undetected failure, which leads to rollback, has probability $p_{\rm fail}=(1-\varpi_p)e^{-T_{\rm par}/T_{\rm det}}$. The value of adopting at $T$ is
\begin{equation}\label{eq:ext-OT}
\mathcal O(T)=\int_T^\infty e^{-rt}\varrho(t-T)\,W(t)\,dt-e^{-rT}\mathcal T(T).
\end{equation}
\begin{enumerate}
\item[(a)] (Base case.) If $T_{\rm par}=0$ and $\mathcal T(T)=F_A(T)=F_0e^{-\delta_FT}$, then $\mathcal O'(T)=e^{-rT}\bigl[(r+\delta_F)F_A(T)-W(T)\bigr]$.
\item[(b)] (Adoption rule.) If $W$ is nondecreasing and $F_A$ nonincreasing, one of them strictly, the net flow $N(T):=W(T)-(r+\delta_F)F_A(T)$ is strictly increasing, and the optimal date is unique: $T^*=0$ if $N(0)\ge0$, otherwise the root of $N$, and never if $N$ stays negative. Adopt when the savings flow first covers the user cost $(r+\delta_F)F_A$ of the integration investment. Under pay-as-you-go pricing, a falling AI cost is a reason to wait only through its effect on today's saving.
\item[(c)] (Option value of waiting.) $\mathcal O(T^*)-\mathcal O(0)=\int_0^{T^*}e^{-rt}[-N(t)]\,dt>0$ whenever $T^*>0$. Adopting today can have a positive net present value, $\mathcal O(0)>0$, while waiting is better, $N(0)<0$.
\item[(d)] (Closed forms.) (i) With a constant savings flow $\bar W$, $T^*=\max\{0,\ln[(r+\delta_F)F_0/\bar W]/\delta_F\}$. (ii) With constant $F_A=F_0$, $v$ and $c_0$, and $c^{\rm var}_A(t)=c^Q+\pi_0e^{-\delta_\pi t}$, where $c^Q$ is the quality cost of Proposition~\ref{prop:tasks-review} (review and expected errors), $T^*=\max\{0,\ln[\pi_0/(c_0-c^Q-rF_0/v)]/\delta_\pi\}$ if $c_0-c^Q>rF_0/v$; otherwise adoption never pays.
\item[(e)] (Vintage lock-in.) If the AI unit cost is fixed at its level on the adoption date, as with a system that cannot be upgraded without a new build, and $v(t)=ve^{g_vt}$ with $g_v<r$, then $\mathcal O'(T)=e^{-rT}\bigl[(r+\delta_F)F_A(T)+v(T)|\dot c^{\rm var}_A(T)|/(r-g_v)-W(T)\bigr]$. The savings flow must then also cover the capitalized decline in AI cost. The ``expected cost decline'' in the hurdle applies only under lock-in.
\item[(f)] (Parallel run and transition costs.) In general
\begin{equation}\label{eq:ext-Opar}
\mathcal O'(T)=e^{-rT}\Bigl[r\,\mathcal T(T)-\mathcal T'(T)-\int_0^{T_{\rm par}}e^{-r\tau}\,W(T+\tau)\,d\varrho(\tau)\Bigr].
\end{equation}
If $W$ is constant over the window, adopt when $\bar\varrho_rW\ge r\mathcal T-\mathcal T'$, where $\bar\varrho_r:=\int_0^{T_{\rm par}}e^{-r\tau}d\varrho(\tau)\in[e^{-rT_{\rm par}},1]$, equal to $(1-e^{-rT_{\rm par}})/(rT_{\rm par})$ for a linear ramp. A component of $\mathcal T$ that falls at rate $\delta_i$ adds $(r+\delta_i)\mathcal T_i$ to the hurdle. A component that grows with volume, as the parallel-run cost does, adds less than $r\mathcal T_i$ and favors earlier adoption. A contract break fee that falls to zero at expiry $T_c$ makes the hurdle drop at $T_c$; if the net flow is negative just before $T_c$ and nonnegative just after, $T^*=T_c$.
\item[(g)] (Capability plateau.) Let capability stop growing at $t_p>0$, so that from $t_p$ on every activity's feasibility and quality cost stay at their values at $t_p$, while inference prices at fixed difficulty may keep falling. Write $W_p\le W$ for the resulting savings flow, equal to $W$ before $t_p$, and $\mathcal O_p$ for the value of adopting.
\begin{enumerate}
\item[(i)] (Pay as you go.) The plateau lowers the value of adopting at every date $T\le t_p-T_{\rm par}$ by the same amount, $\int_{t_p}^\infty e^{-rt}[W(t)-W_p(t)]\,dt$, and at later dates by no more. Where the value of adopting is single-peaked, as under the conditions of (b), the adoption date is weakly later, and in the base case adoption never pays if $W_p(t)<(r+\delta_F)F_A(t)$ at every date. Because adopting now loses at least as much as adopting at any later date, the gain from waiting, $\max_T\mathcal O_p(T)-\mathcal O_p(0)$, does not fall.
\item[(ii)] (Lock-in.) Under the vintage lock-in of (e), the value of adopting at any $T\le t_p$ is unchanged and the value of adopting later falls, so the option value of waiting for better technology, $\max_T\mathcal O_p(T)-\mathcal O_p(0)$, weakly falls; after $t_p$ the capitalized cost decline in the hurdle of (e) keeps only its price part.
\item[(iii)] Activities above the plateau of Proposition~\ref{prop:ext-feasible}(f) never adopt mode A.
\end{enumerate}
\end{enumerate}
\end{proposition}

\noindent\emph{Intuition and example.} The adoption rule is the user-cost rule for an asset whose price is falling. Waiting saves the interest on the build cost and the decline in its price, and costs the savings forgone. With a savings flow of \$4 million a year, $F_0=\$20$ million falling 30\% a year and $r=0.10$, adopting today is worth \$20.0 million, yet the optimal date is 2.31 years away and is worth \$23.8 million. A six-month linear ramp ($\bar\varrho_r=0.975$), \$1.07 million of duplication and reconciliation and \$2 million of stranded costs delay it to 2.67 years. A \$3 million vendor break fee that falls linearly to zero at a contract expiry in year 3 moves it to the expiry date. Migrations bunch at contract renewals for this reason. The general logic is that of \citet{mcdonald1986value} and \citet{dixit1994investment}, and with uncertain arrivals of better technology adoption should be slower still \citep{farzin1998optimal}. Competition can cut the wait short (Proposition~\ref{prop:ext-should}(e)). A plateau in capability cuts both ways, depending on the contract. When AI is paid as used, the firm gets future improvements whenever it adopts, so a plateau lowers the savings of early and late adopters alike, and adoption comes later or not at all. When a system is locked in at its vintage, part of the reason to wait is the better system that later adoption would buy, and a plateau removes it.

\subsection{Competition: should versus could}\label{sec:ext-competition}

\begin{lemma}[Retained savings under competition]\label{lem:ext-retained}
At a date in the market of Assumption~\ref{ass:entry}, let $n_c$ firms with unit costs $\hat c_i$ compete \`a la Cournot with all outputs positive. Then $P=(A+\sum_i\hat c_i)/(n_c+1)$ and firm $i$ earns $(P-\hat c_i)^2/B$ (times $e^{Gt}$). Suppose firm $i$ and $n_r$ of its rivals cut unit cost by $\delta_c$, and outputs stay positive.
\begin{enumerate}
\item[(a)] The price falls by $(n_r+1)\delta_c/(n_c+1)$. Firm $i$'s unit margin rises by $(n_c-n_r)\delta_c/(n_c+1)$, and the unit margin of each firm that does not cut falls by $(n_r+1)\delta_c/(n_c+1)$.
\item[(b)] Firm $i$ keeps the share $1-\bar\varphi=(n_c-n_r)/(n_c+1)$ of its saving, where $\bar\varphi=-dP/d\hat c_i$ is the level pass-through and Lemma~\ref{lem:prices}'s elasticity is $\varphi=\bar\varphi\,u/P$. As $n_c\to\infty$ with $n_r/(n_c-1)\to\mathrm{rep}\in[0,1]$, the retained share tends to $1-\mathrm{rep}$. A saving that no rival replicates is kept, and one that all rivals replicate passes to customers.
\end{enumerate}
\end{lemma}

\begin{assumption}[AI-native entry]\label{ass:ext-entrant}
An incumbent with price $P_I$ and best attainable unit cost $\hat c_I^*$ serves an installed base whose customers switch only if another offer is cheaper by more than a switching premium $p_{\rm sw}$ per unit. AI-native entrants can enter freely, with replication cost tending to zero as in Proposition~\ref{prop:entry}(b), at full unit cost $\hat c_E(t)$. They have no legacy labor and no stranded integration, they have $\iota=1$, and they automate every activity that is feasible, permissible and economical for them.
\end{assumption}

\begin{definition}[Motives for automation]\label{def:ext-motives}
For a feasible and permissible activity $j$ of the incumbent, let $\mathcal B_j$ be its automation advantage per unit (Definition~\ref{def:ext-effcost}), $\mathcal T_j$ its one-time transition cost and $\bar\varphi_j$ the pass-through of its saving in the current regime, and set $\mathrm{NPV}_j:=(1-\bar\varphi_j)\mathcal B_jv_j/(r-G)-\mathcal T_j$. The automation of $j$ is (i) \emph{value-creating} if $\mathrm{NPV}_j>0$; (ii) \emph{defensive} if $\mathrm{NPV}_j\le0$ but $j$ belongs to a least-cost set of activities whose automation keeps the incumbent viable; (iii) \emph{feasible but not worth doing} (could but should not) if it is neither, or if it is feasible but not permissible; and (iv) \emph{not feasible} otherwise.
\end{definition}

\begin{proposition}[Should versus could]\label{prop:ext-should}
\begin{enumerate}
\item[(a)] (Commodity automation among symmetric rivals.) Let all $n_c$ firms, each with unit margin $\mathrm{mg}\ge(n_c-1)\delta_c/(n_c+1)$ so that no firm is driven out, be able to cut unit cost by $\delta_c$ at an annualized fixed cost $\mathcal K$. Let $\mathcal E^{\rm eq}:=[(\mathrm{mg}+\delta_c/(n_c+1))^2-\mathrm{mg}^2]/B$ be a firm's gain when all adopt, and $\mathcal E^{\rm priv}(n_r):=[(\mathrm{mg}+(n_c-n_r)\delta_c/(n_c+1))^2-(\mathrm{mg}-n_r\delta_c/(n_c+1))^2]/B$ its gain from adopting when $n_r$ rivals adopt. Then $\mathcal E^{\rm priv}$ is decreasing in $n_r$, and $\mathcal E^{\rm priv}(n_c-1)\ge\mathcal E^{\rm eq}$. The automation is value-creating if $\mathcal K<\mathcal E^{\rm eq}$. It is defensive if $\mathcal E^{\rm eq}\le\mathcal K<\mathcal E^{\rm priv}(n_c-1)$: adopting is a dominant strategy, yet every firm ends no better off than before, because the saving has passed to customers. It is contested, with an equilibrium in which only some firms adopt, if $\mathcal E^{\rm priv}(n_c-1)\le\mathcal K<\mathcal E^{\rm priv}(0)$, and not worth doing if $\mathcal K\ge\mathcal E^{\rm priv}(0)$. As $n_c\to\infty$, $\mathcal E^{\rm eq}\to0$.
\item[(b)] (Market breakpoint.) Under Assumption~\ref{ass:ext-entrant}, entry binds if and only if $P_I>\hat c_E(t)+p_{\rm sw}$. From the first date $T_M$ at which it does, the incumbent's price is $\hat c_E+p_{\rm sw}$, it keeps all of its own savings ($\bar\varphi=0$, because entrants' cost sets the price), and its unit margin is $\hat c_E+p_{\rm sw}-\hat c_I$.
\item[(c)] (Viability.) After the breakpoint the incumbent survives if and only if $\hat c_I^*\le\hat c_E+p_{\rm sw}$. Write $\hat c_I^*-\hat c_E=\mathcal R_\iota+\mathcal R_{\rm leg}$, the cost of its readiness gap, $\sum_j[c^*_j(\iota^I_j)-c^*_j(1)]$ per unit of output, plus legacy costs. Survival then requires that the switching premium plus any unit-cost advantage from scarce inputs (proprietary data or context, licenses, distribution) cover $\mathcal R_\iota+\mathcal R_{\rm leg}$.
\item[(d)] (Defensive set.) Let $\mathcal D:=\hat c_I^{\rm vc}-(\hat c_E+p_{\rm sw})>0$ be the viability deficit, where $\hat c_I^{\rm vc}$ is the incumbent's unit cost after all value-creating automation. A least-cost defensive set solves $\min_D\sum_{j\in D}(-\mathrm{NPV}_j)$ subject to $\sum_{j\in D}\mathcal B_jv_j/Y_I\ge\mathcal D$, where $Y_I$ is the incumbent's output, over activities with $\mathcal B_j>0$; activities that are not economical never help. The incumbent defends only if this cost is below its franchise value. Ordering activities by $(-\mathrm{NPV}_j)/(\mathcal B_jv_j/Y_I)$ and taking them in turn solves the linear relaxation, with at most one activity taken fractionally.
\item[(e)] (Deadline.) Defensive automation must be complete by $T_M$, so for $j\in D^*$ the adoption date of Proposition~\ref{prop:ext-timing} is replaced by $\min\{T^*_j,\,T_M-\text{lead time}\}$.
\item[(f)] (Capability plateau.) Under a plateau at $t_p$: (i) activities above the plateau of Proposition~\ref{prop:ext-feasible}(f) are not feasible for good, class (iv) of Definition~\ref{def:ext-motives}; (ii) the entrant's cost path rises weakly, $\hat c_{E,p}(t)\ge\hat c_E(t)$ with equality before $t_p$, so the market flip $T_M$ of (b) comes weakly later, and never if $P_I\le\hat c_{E,p}(t)+p_{\rm sw}$ at every date, in which case no automation is defensive; (iii) the viability test of (c) compares the incumbent's best attainable cost with the entrant's, so the plateau eases it only if it raises the entrant's cost by more than the incumbent's; (iv) where the automation advantage $\mathcal B_j$ would have kept growing with capability, the plateau lowers $\mathrm{NPV}_j$, and marginal value-creating activities become feasible but not worth doing.
\end{enumerate}
\end{proposition}

\noindent\emph{Intuition.} What is replicable is competed away. A saving that every rival or entrant can make passes to customers, so automating it is not value-creating, though failing to automate it is ruinous: that is defensive automation. Only savings built on something rivals lack, such as better data, context, process know-how or regulatory standing, raise the firm's long-run margin, consistent with the scarce-input rents $\rho$ of Proposition~\ref{prop:entry}(c). For example, with $\mathrm{mg}=0.20$, $\delta_c=0.10$ and four firms, $\mathcal E^{\rm eq}=0.0084/B$, $\mathcal E^{\rm priv}(3)=0.0288/B$ and $\mathcal E^{\rm priv}(0)=0.0384/B$, so any fixed cost between $0.0084/B$ and $0.0288/B$ makes the automation defensive. For an AI-native entrant whose full unit cost falls as $\hat c_E(t)=0.45+0.75e^{-0.4t}$ against an incumbent price of 1.00 and a switching premium of 0.05, the market flips at $T_M=\ln1.5/0.4=1.01$ years and the long-run price is 0.50. An incumbent whose best attainable cost is 0.55, because of a readiness gap of 0.10, is not viable however much it automates, unless its switching premium or a scarce input is worth at least 0.10 per unit. If capability stops growing after half a year, the entrant's cost stays at 1.06, entry never binds and none of the incumbent's automation is defensive. If it stops after two years, the flip date is unchanged and the long-run price settles at 0.84 instead of 0.50, but an incumbent with the same 0.10 gap is still not viable.

% Slot for the strategic-adoption game of the IO theorist: part2/theory_competition.tex, \label{sec:ext-game}.
% It continues the competition results above as a subsection of this section. Uncomment the next line once the file exists.
\subsection{Strategic adoption among incumbents}\label{sec:ext-game}

Proposition~\ref{prop:ext-timing} dates one firm's adoption against a falling cost, and Proposition~\ref{prop:ext-should}(a) compares outcomes once all rivals have made the same saving; but rivals choose their dates too. We model adoption as a game of timing with the canonical models of \citet{reinganum1981diffusion,reinganum1981market} and \citet{fudenberg1985preemption}, whose timing structure is that of Proposition~\ref{prop:ext-timing}, paired with the modern theory of pass-through \citep{weyl2013pass}, $n$-player preemption \citep{argenziano2014clustering} and the AI-native entrant of Assumption~\ref{ass:ext-entrant}. Proofs are in Appendix~\ref{app:proofs-game}; \texttt{code/verify\_part2\_competition.py} checks every claim and number. Examples are illustrative.

\begin{assumption}[Adoption game]\label{ass:ext-game}
There are $n_c\ge2$ incumbents. Each can adopt once, at a date of its choice, a bundle of automation that lowers its full unit cost by $\delta_c=\sum_j\mathcal B_jv_j/Y>0$ (Definition~\ref{def:ext-effcost}), where $Y$ is its output. Adopting at $T$ costs $F_A(T)=F_0e^{-\delta_FT}$, $F_0>0$, $\delta_F\ge0$, as in Proposition~\ref{prop:ext-timing}(a), and $F_A$ may include the other transition costs $\mathcal T$. Adoptions are observed at once. At date $t$ a firm in state $a\in\{0,1\}$ ($a=1$ once it has adopted), $k$ of whose rivals have adopted, earns $e^{Gt}\mathcal E_a(k)$, with $r>G>-\delta_F$. The private gain $\mathcal E^{\rm priv}(k):=\mathcal E_1(k)-\mathcal E_0(k)$ is positive and strictly decreasing in $k$ (adoptions are strategic substitutes), and $\mathcal E_1$ and $\mathcal E_0$ are strictly decreasing in $k$ (a rival's adoption hurts). The gain when all adopt is $\mathcal E^{\rm eq}:=\mathcal E_1(n_c-1)-\mathcal E_0(0)$, as in Proposition~\ref{prop:ext-should}(a). The single-firm date of a gain flow $\mathcal E$,
\begin{equation}\label{eq:ext-game-To}
T^\circ(\mathcal E):=\max\Bigl\{0,\ \frac1{G+\delta_F}\ln\frac{(r+\delta_F)F_0}{\mathcal E}\Bigr\}\ \ (\mathcal E>0),\qquad T^\circ(\mathcal E):=\infty\ \ (\mathcal E\le0),
\end{equation}
is the date at which $e^{Gt}\mathcal E$ first covers the user cost $(r+\delta_F)F_A(t)$ (Proposition~\ref{prop:ext-timing}(b)).
\end{assumption}

With two firms, let $T_2:=T^\circ(\mathcal E^{\rm priv}(1))$. Following \citet{fudenberg1985preemption}, write $V^{\rm L}(t)$ for the value at date 0 of a firm that adopts first at $t$ while its rival follows at its best date $\max\{t,T_2\}$, $V^{\rm F}(t)$ for the rival's value, and $V^{\rm J}(t)$ for each firm's value when both adopt at $t$. With $I(a,b):=\int_a^be^{-(r-G)s}ds$, for $t\le T_2$,
\begin{equation}\label{eq:ext-game-values}
\begin{aligned}
V^{\rm L}(t)&=\mathcal E_0(0)I(0,t)+\mathcal E_1(0)I(t,T_2)+\mathcal E_1(1)I(T_2,\infty)-e^{-rt}F_A(t),\\
V^{\rm F}(t)&=\mathcal E_0(0)I(0,t)+\mathcal E_0(1)I(t,T_2)+\mathcal E_1(1)I(T_2,\infty)-e^{-rT_2}F_A(T_2),\\
V^{\rm J}(t)&=\mathcal E_0(0)I(0,t)+\mathcal E_1(1)I(t,\infty)-e^{-rt}F_A(t),
\end{aligned}
\end{equation}
and $V^{\rm L}=V^{\rm F}=V^{\rm J}$ for $t\ge T_2$.

\begin{lemma}[Pass-through and the adoption incentives]\label{lem:ext-game-pass}
\begin{enumerate}
\item[(a)] (Linear demand.) Let $\bar\varphi^{\rm own}$ and $\bar\varphi^{\rm ind}$ be the level pass-through of a saving that no rival makes and of one that all make (Lemma~\ref{lem:ext-retained}(b)), and $\bar\varphi^\times:=(\bar\varphi^{\rm ind}-\bar\varphi^{\rm own})/(n_c-1)$ the effect of one rival's saving on a firm's price. If, with unit margin mg before adoption, $\mathcal E_1(k)=[\mathrm{mg}+(1-\bar\varphi^{\rm own})\delta_c-k\bar\varphi^\times\delta_c]^2/B_q$ and $\mathcal E_0(k)=[\mathrm{mg}-k\bar\varphi^\times\delta_c]^2/B_q$, with $\bar\varphi^\times>0$ and $\mathrm{mg}>(n_c-1)\bar\varphi^\times\delta_c$, then the conditions of Assumption~\ref{ass:ext-game} on $\mathcal E_a$ hold. The Cournot market of Lemma~\ref{lem:ext-retained} has this form with $B_q=B$, $\bar\varphi^{\rm own}=\bar\varphi^\times=1/(n_c+1)$ and $\bar\varphi^{\rm ind}=n_c/(n_c+1)$; so does differentiated Bertrand competition with demand $q_i=A_B-p_i+s_B\sum_{l\ne i}p_l$, $0<s_B<1/(n_c-1)$, with $B_q=1$, $\bar\varphi^{\rm ind}=1/[2-(n_c-1)s_B]$ and $\bar\varphi^\times=s_B\bar\varphi^{\rm ind}/(2+s_B)$.
\item[(b)] (Cournot with general demand.) Let inverse demand $P(Q)$ be twice continuously differentiable with $P'<0$, all firms have unit cost $\hat c$ and output $Y$, $\mathsf e:=-QP''(Q)/P'(Q)$ be the curvature of inverse demand, and the symmetric equilibrium be regular, $n_c+1-\mathsf e>0$. Then $\bar\varphi^{\rm ind}=n_c/(n_c+1-\mathsf e)$, $\bar\varphi^{\rm own}=\bar\varphi^{\rm ind}/n_c$, Lemma~\ref{lem:prices}'s elasticity is $\varphi=\bar\varphi^{\rm ind}\hat c/P$, and as $\delta_c\to0$, for every $k$,
\begin{equation}\label{eq:ext-game-first}
\mathcal E^{\rm eq}=Y\delta_c\Bigl[1-\bar\varphi^{\rm ind}\Bigl(1-\frac1{n_c}\Bigr)\Bigr]+o(\delta_c),\qquad \mathcal E^{\rm priv}(k)=Y\delta_c\Bigl[\frac{2n_c-1}{n_c}-\frac{(n_c-1)\bar\varphi^{\rm ind}}{n_c^2}\Bigr]+o(\delta_c).
\end{equation}
The first is the incidence formula of \citet{weyl2013pass} with Cournot conduct $1/n_c$.
\item[(c)] (Constant elasticity.) If $Q=DP^{-\epsilon}$ with $n_c\epsilon>1$, then $P=n_c\epsilon\hat c/(n_c\epsilon-1)$, $\bar\varphi^{\rm ind}=n_c\epsilon/(n_c\epsilon-1)$, $\varphi=1$ and the Lerner index is $1/(n_c\epsilon)$. Earnings and revenue are proportional to $\hat c^{\,1-\epsilon}$, so $\mathcal E^{\rm eq}=\mathcal E_0(0)[(1-\delta_c/\hat c)^{1-\epsilon}-1]$ has the sign of $\epsilon-1$. The private gain in \eqref{eq:ext-game-first} is positive if and only if $\epsilon>(2n_c-1)/(2n_c^2-2n_c+1)$.
\end{enumerate}
\end{lemma}

\noindent\emph{Intuition and example.} Pass-through decides how much of a saving leaves the industry; the volume response decides how much a firm gains by taking business from rivals. With four Cournot firms the two pass-throughs are 20\% and 80\%, and $\varphi=0.64$ at a price of 1.00 and unit cost 0.80; under Bertrand competition with $s_B=0.2$ they are 0.519 and 0.714. Under constant elasticity $\epsilon=0.8$ with four firms, $\bar\varphi^{\rm ind}=1.455$: a common saving of 0.10 on a unit cost of 0.6875 cuts the price by 14.5\% and every firm's earnings by 3.1\%, though adopting alone still pays (the threshold in (c) is 0.28 with four firms, 0.6 with two).

\begin{proposition}[Diffusion and preemption]\label{prop:ext-game-preempt}
Under Assumption~\ref{ass:ext-game}:
\begin{enumerate}
\item[(a)] (Precommitment.) Let $T_k:=T^\circ(\mathcal E^{\rm priv}(k-1))$, $k=1,\dots,n_c$. If firms commit to dates at date 0, the dates $T_1\le\dots\le T_{n_c}$, assigned in any order, form a Nash equilibrium. They increase strictly where positive, so identical firms adopt in sequence \citep{reinganum1981diffusion,reinganum1981market}. The $k$-th adopter moves before $T^\circ(\delta_cY)$, the date of a firm that takes prices and volume as given, if and only if $\mathcal E^{\rm priv}(k-1)>\delta_cY$.
\item[(b)] (Preemption.) Let $n_c=2$, $T_2>0$ and $V^{\rm L}(0)<V^{\rm F}(0)$, and let $T_1:=T^\circ(\mathcal E^{\rm priv}(0))$ and $T_D:=T^\circ(\mathcal E_1(0)-\mathcal E_0(1))$. Then $0<T_D<T_1<T_2$, and $V^{\rm L}-V^{\rm F}$ rises strictly on $[0,T_D]$, falls strictly on $[T_D,T_2]$ and vanishes at $T_2$, so it has a unique root $T^{\rm P}\in(0,T_D)$. If firms can respond to each other, the timing game of \citet{fudenberg1985preemption} has a subgame-perfect equilibrium in which one firm adopts at $T^{\rm P}$ and the other at $T_2$. Rents are equalized at $V^{\rm L}(T^{\rm P})=V^{\rm F}(T^{\rm P})$, and the race costs the leader $V^{\rm L}(T_1)-V^{\rm L}(T^{\rm P})>0$.
\end{enumerate}
\end{proposition}

\noindent\emph{Intuition and example.} A committed firm adopts when its private gain covers the user cost, and since each adoption lowers the next firm's gain, the dates fan out. A firm that can be preempted cannot wait for its best date: while leading is worth more than following, a rival would move first, so the first adoption is pulled forward until the two are worth the same. Take four Cournot incumbents, each selling 20 million units a year at \$1.00 with unit cost \$0.80 and able to save \$0.10 a unit, with $r=0.10$, $G=0.04$ and, as in Proposition~\ref{prop:ext-timing}'s example, an adoption cost of \$20 million falling 30\% a year. The private gains are \$3.84, 3.52, 3.20 and 2.88 million a year (Proposition~\ref{prop:ext-should}(a) with $B=0.01$), and the precommitment dates 2.16, 2.41, 2.69 and 3.00 years, all before the 4.08 years at which a firm ignoring its rivals would adopt the \$2.0 million engineering saving: the private gain includes business taken from rivals. The durable gain of \$0.84 million would cover the user cost only after 6.63 years. In a duopoly with the same numbers per firm, $(\mathcal E_0(0),\mathcal E_1(0),\mathcal E_0(1),\mathcal E_1(1))=(4.00,7.11,2.78,5.44)$, $T_D=1.80$, $T_1=2.78$ and $T_2=3.23$, and preemption brings the first adoption forward to $T^{\rm P}=0.63$ years; each firm is worth \$78.2 million, against \$81.1 million for an unchallenged leader at $T_1$ and \$66.7 million if no one adopted. For $n_c>2$ we report the precommitment dates, because rent equalization does not extend to oligopoly \citep{fudenberg1985preemption} and $n$-player preemption can cluster investments \citep{argenziano2014clustering}.

\begin{proposition}[Waiting games]\label{prop:ext-game-wait}
Under Assumption~\ref{ass:ext-game} with $n_c=2$ and $T_2>0$:
\begin{enumerate}
\item[(a)] (Late joint adoption.) Let $T_J:=T^\circ(\mathcal E^{\rm eq})$, which exceeds $T_2$ and is infinite if $\mathcal E^{\rm eq}\le0$. For finite $T\in[T_2,T_J]$, joint adoption at $T$ is a subgame-perfect equilibrium if and only if $V^{\rm J}(T)\ge V^{\rm L}(T_1)$, and no joint date outside $[T_2,T_J]$ is. Such equilibria exist (never adopting, if $T_J=\infty$) if and only if
\begin{equation}\label{eq:ext-game-joint}
\int_{T_2}^{T_J}e^{-rt}\bigl[(r+\delta_F)F_A(t)-\mathcal E^{\rm eq}e^{Gt}\bigr]dt\ \ge\ \int_{T_1}^{T_2}e^{-rt}\bigl[\mathcal E^{\rm priv}(0)e^{Gt}-(r+\delta_F)F_A(t)\bigr]dt,
\end{equation}
and they give each firm more than the preemption equilibrium.
\item[(b)] (Learning.) Let the bundle work for the industry with known probability $\varpi\in(0,1)$, the first adoption reveal at once whether it does, and a failure leave the adopter's cost unchanged, its outlay sunk and no further adoption. Then $V^{\rm L}_\varpi(t):=\varpi V^{\rm L}(t)+(1-\varpi)[\mathcal E_0(0)/(r-G)-e^{-rt}F_A(t)]$ and $V^{\rm F}_\varpi(t):=\varpi V^{\rm F}(t)+(1-\varpi)\mathcal E_0(0)/(r-G)$. If $\varpi\le\bar\varpi:=\mathcal E^{\rm priv}(1)/[\mathcal E_1(0)-\mathcal E_0(1)]\in(0,1)$, then $V^{\rm F}_\varpi>V^{\rm L}_\varpi$ at every date: no firm ever gains by preempting.
\item[(c)] (War of attrition.) Under (b), let $\varpi\le\bar\varpi$, $\mathcal E^{\rm eq}>0$ and $T^{\rm W}:=T^\circ(\varpi\mathcal E^{\rm eq})\ge T_2$. Then $V^{\rm L}_\varpi$ rises strictly before $T^{\rm W}$ and falls strictly after it. (i) For either assignment of roles there is an equilibrium in which the leader adopts at $T^{\rm W}$ and the follower adopts at once if the bundle works and never otherwise. (ii) There is a symmetric equilibrium in which, until someone adopts, each firm adopts at the hazard rate
\begin{equation}\label{eq:ext-game-hazard}
\nu^{\rm W}(t)=\frac{\varpi\mathcal E^{\rm eq}e^{Gt}-(r+\delta_F)F_A(t)}{(1-\varpi)F_A(t)}\quad(t\ge T^{\rm W}),\qquad \nu^{\rm W}(t)=0\quad(t<T^{\rm W});
\end{equation}
each firm is worth $V^{\rm L}_\varpi(T^{\rm W})$, less than the follower in (i), and the first adoption comes a random delay $s$ after $T^{\rm W}$, with $\Pr(s>u)=\exp\{-2\int_{T^{\rm W}}^{T^{\rm W}+u}\nu^{\rm W}\}$.
\end{enumerate}
\end{proposition}

\noindent\emph{Intuition and example.} Two forces make firms wait. If the gain from leading is small, both can wait for the late date at which the durable gain covers the user cost, since a firm that broke ranks would gain less from leading than it loses by triggering early adoption \citep{fudenberg1985preemption}. If success is uncertain, the first adopter reveals the answer to its rival, who copies a success and avoids a failure; following beats leading, and adoption becomes a war of attrition \citep{hendricks1988war} with second-mover advantages \citep{hoppe2000second}, timed by the risk-adjusted durable gain $\varpi\mathcal E^{\rm eq}$. In the duopoly above, joint adoption at any date in $[3.32,5.03]$ is an equilibrium worth up to \$81.8 million to each firm, against \$78.2 million under preemption. Such equilibria exist here at every cut for which the follower waits ($T_2>0$); they fail for large cuts and costs, and with a cut of \$0.31 and an adoption cost of \$80 million the outcome is preemption at 0.66 years. If the bundle works with probability $\varpi=0.5$, below $\bar\varpi=0.615$ (preemption reappears only above 0.893), the war of attrition starts at $T^{\rm W}=7.07$ years, and the first adoption follows after a mean delay of 1.52 years (median 1.46); each firm is worth \$73.4 million, against \$74.0 million for an assigned follower.

\begin{proposition}[The automation prisoner's dilemma]\label{prop:ext-game-pd}
\begin{enumerate}
\item[(a)] (Static.) Let $n_c$ symmetric firms each decide whether to adopt at an annualized cost $\mathcal K$, with $\mathcal E^{\rm priv}(k)$ decreasing and $\mathcal E_0(k)$ strictly decreasing in $k$. Adopting is strictly dominant if and only if $\mathcal K<\mathcal E^{\rm priv}(n_c-1)$, and when all adopt each firm is no better off if and only if $\mathcal K\ge\mathcal E^{\rm eq}$. The dilemma window $[\mathcal E^{\rm eq},\mathcal E^{\rm priv}(n_c-1))$ has width $\mathcal E_0(0)-\mathcal E_0(n_c-1)>0$; to first order in symmetric Cournot it is
\begin{equation}\label{eq:ext-game-window}
\frac{\mathcal K}{Y\delta_c}\in\Bigl[1-\bar\varphi^{\rm ind}\Bigl(1-\frac1{n_c}\Bigr),\ \frac{2n_c-1}{n_c}-\frac{(n_c-1)\bar\varphi^{\rm ind}}{n_c^2}\Bigr),
\end{equation}
which under linear demand is $[2/(n_c+1),2n_c/(n_c+1))$: only the share $1/n_c$ of the private gain is durable.
\item[(b)] (Pass-through, elasticity and the number of firms.) To first order, common adoption lowers every firm's gross earnings, so that the dilemma holds at every $\mathcal K\ge0$ at which adopting is privately worthwhile, if and only if $\bar\varphi^{\rm ind}>n_c/(n_c-1)$, that is $\mathsf e>2$, whatever $n_c$. Under constant elasticity it holds exactly if and only if $\epsilon<1$; under linear demand, never.
\item[(c)] (Dynamic.) Under Assumption~\ref{ass:ext-game}, in the precommitment equilibrium the last adopter is no better off than if no firm ever adopted whenever $\mathcal E^{\rm eq}\le(r-G)e^{-GT_{n_c}}F_A(T_{n_c})$, which for $T_{n_c}>0$ reads
\begin{equation}\label{eq:ext-game-dyn}
\mathcal E^{\rm eq}\le\frac{r-G}{r+\delta_F}\,\mathcal E^{\rm priv}(n_c-1).
\end{equation}
Under the conditions of Proposition~\ref{prop:ext-game-preempt}(b), the same condition with $n_c=2$ leaves both duopolists no better off in the preemption equilibrium. Under linear demand it holds, to first order, when $n_c\ge(r+\delta_F)/(r-G)$.
\item[(d)] (Rule of $X$.) In symmetric Cournot with $b=0$, let $S_0$, $S_1$ be a firm's revenue before and after common adoption and $m_0$ its margin before. Its margin net of $\mathcal K$ changes by exactly $[\mathcal E^{\rm eq}-\mathcal K-m_0(S_1-S_0)]/S_1$. To first order, revenue changes by $-(1-\epsilon)\varphi\delta_c/\hat c$ (recorded once in growth, with weight $X$) and the margin before $\mathcal K$ by $(1-\varphi)\delta_c/P$, and a continuing fall in unit cost at rate $\delta_u$ lowers nominal growth by $(1-\epsilon)\varphi\delta_u$ (Lemma~\ref{lem:prices}). Under constant elasticity the margin falls by exactly $\mathcal K/S_1$; if also $\epsilon<1$, every incumbent's Rule-of-$X$ score falls in both terms, and its value $Mv(s)$ (Proposition~\ref{prop:dynamic}(a)) falls twice over at unchanged shares, as $M$ shrinks by the factor $(1-\delta_c/\hat c)^{1-\epsilon}$ and $v$ falls with the margin intercept.
\end{enumerate}
\end{proposition}

\noindent\emph{Intuition and example.} Each firm adopts because a saving it alone made would take business from its rivals. When all adopt, no business moves and the price falls, leaving the durable share: small when pass-through is high and firms are many, negative when pass-through exceeds $n_c/(n_c-1)$, because industry earnings then fall with industry costs \citep{weyl2013pass}. In the four-firm example the window is $[0.84,2.88)$ million dollars a year, against $[0.8,3.2)$ to first order; under constant elasticity $\epsilon=0.8$, common adoption costs every firm 3.1\% of earnings even if free. The dynamic test compares the durable gain with the share $(r-G)/(r+\delta_F)=0.15$ of the last adopter's private gain, the annualized cost of its investment at its chosen date. With four firms the ratio is 0.29 and the last adopter gains; with ten it is 0.125 and the last adopter loses \$2.73 million (the first-order bound is $n_c\ge6.67$). Falling adoption costs soften the dilemma: with a constant cost of \$20 million all four firms adopt at once, and each loses \$6.0 million although the last adopter's private gain is worth \$48 million. On the Rule of 40, common adoption in the linear market raises revenue by 1.2\% and the margin before adoption cost from 20.0\% to 23.9\%, but at $\mathcal K=\mathcal E^{\rm eq}$ the margin is 19.8\%; under constant elasticity $\epsilon=0.8$, revenue falls by 3.1\% with the margin of 31.25\% unchanged, and a continuing 5\% annual fall in unit cost costs 1 point of growth. Faster growth of AI-investing firms \citep{babina2024artificial} fits the business-stealing part of the private gain, which is not durable for the industry.

\begin{proposition}[AI-native entry and defensive automation]\label{prop:ext-game-entry}
Let the incumbents be the Cournot firms of Lemma~\ref{lem:ext-game-pass}(a), with price $P_0$ and unit cost $\hat c_0$ before AI and $P_k=P_0-k\delta_c/(n_c+1)$ once $k$ have adopted, and the entrants as in Assumption~\ref{ass:ext-entrant}, with full unit cost $\hat c_E(t)$ strictly decreasing to $\hat c_E^\infty$; write $P_E(t):=\hat c_E(t)+p_{\rm sw}$ and $P_E^\infty:=\hat c_E^\infty+p_{\rm sw}$. Entry binds when the incumbents' price exceeds $P_E(t)$ (Proposition~\ref{prop:ext-should}(b)); each incumbent then sells its installed base $Ye^{Gt}$ at $P_E(t)$.
\begin{enumerate}
\item[(a)] (Entry ends the game.) From $\bar T_M:=\inf\{t:P_E(t)<P_{n_c}\}$ on, entry binds whatever incumbents do, and an incumbent's gain from adopting is $\delta_cYe^{Gt}$ whatever its rivals do. Pass-through is zero, the motives of Propositions~\ref{prop:ext-game-preempt} and \ref{prop:ext-game-wait}(a) disappear (the learning externality of Proposition~\ref{prop:ext-game-wait}(b) remains), and an incumbent that stays and has not adopted at $t\ge\bar T_M$ adopts at $\max\{t,T^\circ(\delta_cY)\}$.
\item[(b)] (Deterrence is a public good.) The breakpoint $T_M(k):=\inf\{t:P_E(t)<P_k\}$ is nondecreasing in the number of adopters $k$, so each adoption postpones entry for all incumbents. If $P_{n_c}\le P_E^\infty$, entry never binds if and only if, for each $j\le\bar k:=\min\{k:P_k\le P_E^\infty\}$, the $j$-th adoption takes place by $T_M(j-1)$.
\item[(c)] (Defensive automation.) Let $P_E^\infty<P_{n_c}$ and let adoption cost $\mathcal K$ a year. In the long run an adopter's flow exceeds a non-adopter's by $\delta_cY-\mathcal K$, since it keeps all of its saving, and exceeds its own flow before AI by $(P_E^\infty-P_0+\delta_c)Y-\mathcal K$. Adoption is therefore defensive, in the sense of Proposition~\ref{prop:ext-should}(a), if and only if $(P_E^\infty-P_0+\delta_c)Y\le\mathcal K<\delta_cY$, and if $P_E^\infty\le P_0-\delta_c$ every adoption worth making is defensive. It is needed for viability if and only if $P_E^\infty<\hat c_0$, and it secures viability if and only if $P_E^\infty\ge\hat c_0-\delta_c+\mathcal K/Y$ (Proposition~\ref{prop:ext-should}(c)).
\item[(d)] (Free replication.) If entrants make the incumbent's saving at least as cheaply, $\hat c_E^\infty\le\hat c_0-\delta_c$, as under the free replication of Proposition~\ref{prop:entry}(b), the incumbent's long-run unit margin is at most $p_{\rm sw}$ plus any cost advantage from a scarce input (Proposition~\ref{prop:entry}(c)).
\end{enumerate}
\end{proposition}

\noindent\emph{Intuition and example.} Once entry binds, entrants' cost sets the price: an incumbent keeps its whole saving but takes no business from rivals, so the strategic motives vanish. The entrant also moves the benchmark: relative to the world before AI, automation no longer raises the incumbent's margin; it only limits the fall. Let entrants' cost be $\hat c_E(t)=0.70+0.50e^{-0.4t}$, with $p_{\rm sw}=0.05$, against the four incumbents. Entry binds at 1.73, 1.94, 2.17, 2.42 or 2.70 years if zero to four incumbents have adopted, and from 2.70 years whatever they do; no schedule deters it, since $P_4=0.92>P_E^\infty=0.75$. After 2.70 years the gain from adopting is the \$2.0 million engineering saving, so a firm that has not adopted waits until 4.08 years, before its margin turns negative at 5.76 years. In the long run an adopter earns a unit margin of 0.05, against 0.20 before AI and $-0.05$ without automation; since $P_E^\infty=0.75\le P_0-\delta_c=0.90$, every adoption worth making is defensive (window $[-3.0,2.0)$ million dollars a year). Proposition~\ref{prop:prod-cannibal}(c) gives the analogous product choice.

\begin{corollary}[Durable gain, concentration and should versus could]\label{cor:ext-game-engine}
Let HHI be the incumbents' Herfindahl index, with $n_c=1/\mathrm{HHI}$ for symmetric firms, under the conditions of Lemma~\ref{lem:ext-game-pass}(b).
\begin{enumerate}
\item[(a)] (Durable share.) To first order a firm keeps the share
\begin{equation}\label{eq:ext-game-dur}
\mathrm{dur}:=1-\bar\varphi^{\rm ind}(1-\mathrm{HHI})=\frac{2-\mathsf e}{1/\mathrm{HHI}+1-\mathsf e}
\end{equation}
of a saving $\delta_cY$ that every rival makes, measured in dollars of profit with its volume gain counted (it is not a margin share: at the lower price the margin keeps $1-\bar\varphi^{\rm ind}u/p$ of the saving to first order, which does not depend on HHI; Proposition~\ref{prop:tasks-aggregate}(a)), and $\mathrm{priv}:=2-\mathrm{HHI}-\mathrm{HHI}(1-\mathrm{HHI})\bar\varphi^{\rm ind}$ of one that none makes ($2\,\mathrm{HHI}/(1+\mathrm{HHI})$ and $2/(1+\mathrm{HHI})$ under linear demand; $\mathrm{dur}=\mathrm{HHI}(\epsilon-1)/(\epsilon-\mathrm{HHI})$ under constant elasticity), and $(1-\mathrm{rep})\,\mathrm{priv}+\mathrm{rep}\cdot\mathrm{dur}$ if the share rep of rivals can replicate it. With constant curvature, dur is one under monopoly and rises with concentration if and only if $\mathsf e<2$. Where entry binds (Proposition~\ref{prop:ext-game-entry}(c)), $\mathrm{dur}=(P_E^\infty-P_0+\delta_c)/\delta_c$ and $\mathrm{priv}=1$.
\item[(b)] (Classification.) For a bundle with advantage $\sum_j\mathcal B_jv_j=\delta_cY$, success probability $\varpi$ and transition cost $\mathcal T$, let $\mathrm{NPV}^{\rm dur}:=\varpi\,\mathrm{dur}\,\delta_cY/(r-G)-\mathcal T$ and $\mathrm{NPV}^{\rm priv}:=\varpi\,\mathrm{priv}\,\delta_cY/(r-G)-\mathcal T$. The bundle is value-creating if $\mathrm{NPV}^{\rm dur}>0$; defensive if $\mathrm{NPV}^{\rm dur}\le0<\mathrm{NPV}^{\rm priv}$ (the strategic case of Proposition~\ref{prop:ext-should}(a)) or if it is needed for viability (Definition~\ref{def:ext-motives}(ii)); and feasible but not worth doing otherwise. Its date is the strategic date of Propositions~\ref{prop:ext-game-preempt} and \ref{prop:ext-game-wait}, capped by the deadline of Proposition~\ref{prop:ext-should}(e).
\item[(c)] (Expected durable gain.) The expected durable gain is $\varpi\,\mathrm{dur}\,\delta_cY-\mathcal K$ a year; when $\mathsf e<2$ it rises with concentration, from $-\mathcal K$ in a fragmented industry to $\varpi\delta_cY-\mathcal K$ under monopoly.
\end{enumerate}
\end{corollary}

\begin{table}[t]
\centering
\caption{Durable and private shares of a replicable saving, to first order (Corollary~\ref{cor:ext-game-engine}(a))}\label{tab:ext-game-dur}
\footnotesize
\begin{tabular}{@{}lcccccc@{}}
\toprule
Number of incumbents $n_c$ (HHI $=1/n_c$) & 1 & 2 & 3 & 4 & 6 & 10\\
\midrule
Linear demand: durable share $2/(n_c+1)$ & 1.00 & 0.67 & 0.50 & 0.40 & 0.29 & 0.18\\
Linear demand: private share $2n_c/(n_c+1)$ & 1.00 & 1.33 & 1.50 & 1.60 & 1.71 & 1.82\\
Constant elasticity $\epsilon=2$: durable share & 1.00 & 0.33 & 0.20 & 0.14 & 0.09 & 0.05\\
Constant elasticity $\epsilon=0.8$: durable share & n.a. & $-$0.33 & $-$0.14 & $-$0.09 & $-$0.05 & $-$0.03\\
\bottomrule
\end{tabular}

\smallskip
\begin{minipage}{0.92\linewidth}\footnotesize
Notes: shares of the engineering saving $\delta_cY$. A monopoly with $\epsilon<1$ has no finite price (n.a.). Values computed by \texttt{code/verify\_part2\_competition.py}.
\end{minipage}
\end{table}

\noindent\emph{Intuition and example.} The engine's question is not whether a saving can be made but how much of it survives the rivals' response. A monopolist keeps all of it; four Cournot firms with linear demand keep 40\% once all adopt, though each would gain 160\% by adopting alone; under inelastic demand the share is negative (Table~\ref{tab:ext-game-dur}). A bundle that saves the example's \$2.0 million a year, works with probability $\varpi=0.6$ and costs $\mathcal T=\$20$ million has $\mathrm{NPV}^{\rm dur}=-\$12.0$ million and $\mathrm{NPV}^{\rm priv}=\$12.0$ million: it is defensive, to be done only because rivals will do it. Definition~\ref{def:ext-motives}, which counts the retained share $1/(n_c+1)$ of the unit margin at fixed volume, gives $-\$16.0$ million; the volume response doubles the durable share under linear demand. For should versus could, three implications follow. In fragmented markets most replicable automation is defensive, worth doing to avoid a loss rather than to raise the margin. In concentrated markets more of it is durable, but duopolies may delay it, by tacit joint waiting or, when success is uncertain, by a war of attrition. Where AI-native entry binds, the durable gain is set by the entrants' cost and the switching premium, and the incumbents' game no longer matters. The results assume symmetric incumbents, one bundle per firm and deterministic adoption costs; with asymmetric firms the engine uses $n_c=1/\mathrm{HHI}$ as an approximation.

\noindent\emph{Calibration of the industry pass-through.} Lemma~\ref{lem:ext-game-pass}(b) derives $\bar\varphi^{\rm ind}=n_c/(n_c+1-\mathsf e)$ for homogeneous Cournot firms, which is 0.8 for four firms under linear demand and at least 0.75 for three or more. That is a benchmark of the model, not a measurement, and it lies above the evidence on cost pass-through. Developers of differentiated apps cut prices on 9\% of products after a cost cut of about ten points of revenue, with a median cut of 2.5\%, within three months (a study funded by Apple, without a control group) \citep{choi2025dma}; about 70\% of industry-wide energy cost changes pass through to U.S.\ manufacturing prices in the short to medium run \citep{ganapati2020energy}; and final pass-through of value-added tax is about 100\% for standard-rate changes and about 30\% for reduced-rate changes in the euro area \citep{benedek2015vat}. Part~II therefore treats $\bar\varphi^{\rm ind}$ as a scenario dimension with central value 0.5, the medium-run case between the software and manufacturing anchors, and scenarios 0 (fixed prices, the short run for differentiated software) and 0.7 (an industry-wide cost shock); values near one belong to undifferentiated services in the long run (all class b). Equation~\eqref{eq:ext-game-dur} is kept with the measured value in place of the derived one; within the lemma's model, 0.5 and 0.7 at four firms correspond to curvatures $\mathsf e=n_c+1-n_c/\bar\varphi^{\rm ind}$ of $-3$ and $-0.71$, that is, to more concave demand, and they stand for the differentiation and frictions the homogeneous model leaves out. At four symmetric firms ($\mathrm{HHI}=0.25$) the durable share is 1.0, 0.625 and 0.475 at pass-through 0, 0.5 and 0.7, against 0.40 at the Cournot benchmark, and the private share is 1.75, 1.66 and 1.62 against 1.60. The bundle of the example then has $\mathrm{NPV}^{\rm dur}$ of \$0, $-\$7.5$ million and $-\$10.5$ million and $\mathrm{NPV}^{\rm priv}$ of \$15.0, \$13.1 and \$12.4 million: it stays defensive at every scenario, and at fixed prices it is on the boundary, worth exactly its cost. These are dollar shares. The rate $\bar\varphi^{\rm ind}$ also sets the margin, and so the score: at the lower price a replicated saving keeps $1-\bar\varphi^{\rm ind}u/p$ of its margin gain to first order, whatever the concentration (Part~I's mapping, which applies the rate to every unit-cost saving and computes the margin at the post-cut price), so a lower pass-through raises the rank of cost-of-revenue automation among the levers of Proposition~\ref{prop:playbook-rank} through that share, not through the durable share; the tables above remain the linear-demand benchmark.


\subsection{Organization: teams, hierarchies and the residual firm}\label{sec:ext-org}

\begin{assumption}[Team production with agents]\label{ass:ext-team}
A team must deliver a scope of $\mathcal W_0$ person-years of unaided work a year. It has $n>0$ humans (fractions are part-time work), each with AI leverage $A_g=t_j/t^U_j$ (the augmentation gain of Section~\ref{sec:tasks}), and for $n\ge1$ the humans lose $\chi_hn(n-1)/2$ person-years a year to coordination. Agents can perform the share $\bar\nu_a$ of the scope; the team assigns them a share $\nu_a\in[0,\bar\nu_a]$, and their work needs $\xi^t$ person-years of human review per person-year of output, which is $\xi_jt^R_j/t_j$ in the notation of Proposition~\ref{prop:tasks-review}. Each of the $n_a$ agents costs $p_a$ a year and produces $Y_a$ person-year equivalents, and agents lose $\chi_an_a(n_a-1)/2$ of output to coordination among themselves. The wage is $w$. The team needs $n-\chi_hn(n-1)_+/2\ge \mathcal W(\nu_a):=\mathcal W_0[(1-\nu_a)/A_g+\nu_a\xi^t]$.
\end{assumption}

The pairwise form of coordination cost is an assumption. It is motivated by the argument of \citet{becker1992division} that coordination costs, rather than the extent of the market, limit the division of labor, and by the observation of \citet{brooks1975mythical} that adding people to a late project makes it later.

\begin{proposition}[Team size and the solo contributor]\label{prop:ext-team}
Let $\chi_a=0$ in (a) to (f).
\begin{enumerate}
\item[(a)] (Agent share.) If $p_a/Y_a<w(1/A_g-\xi^t)$, which is the activity-level condition $c_{\mathrm A}<c_{\mathrm U}$, agents take all the work they can do, $\nu_a=\bar\nu_a$. Otherwise the cost-minimizing share is zero or interior, where the coordination-adjusted wage $\tilde w=w/[1-\chi_h(n-\tfrac12)]$ satisfies $\tilde w(1/A_g-\xi^t)=p_a/Y_a$. Coordination costs therefore make automation pay in large teams where it does not pay activity by activity.
\item[(b)] (Team size.) The minimal team is $n^*(\mathcal W)=\mathcal W$ for $\mathcal W\le1$ and
\begin{equation}\label{eq:ext-nstar}
n^*(\mathcal W)=\frac{(1+\chi_h/2)-\sqrt{(1+\chi_h/2)^2-2\chi_h\mathcal W}}{\chi_h}\qquad\text{for }1<\mathcal W\le \mathcal W_{\max}:=\frac{(1+\chi_h/2)^2}{2\chi_h}.
\end{equation}
A flat team cannot deliver more than $\mathcal W_{\max}$; beyond $n=1/\chi_h+1/2$ an extra member lowers output.
\item[(c)] (Solo breakpoint.) $n^*\le1$ if and only if $\mathcal W_0\le\Lambda:=1/[(1-\nu_a)/A_g+\nu_a\xi^t]$, the scope a single contributor directing agents can deliver.
\item[(d)] (Amdahl bounds.) $\Lambda\le A_g/(1-\nu_a)$ and $\Lambda\le1/(\nu_a\xi^t)$, and $\Lambda\to1/\xi^t$ as $\nu_a\to1$. An oversight floor caps leverage at $1/(\nu_a\xi^{\min}t^R/t)$. A tenfold contributor needs $(1-\nu_a)/A_g+\nu_a\xi^t\le0.1$; with $A_g\le2$ this requires $\nu_a\ge0.8$ and hence $\xi^t<0.125$.
\item[(e)] (Coordination dividend.) $dn^*/d\mathcal W=1/[1-\chi_h(n^*-\tfrac12)]\ge1$, and $n^*$ is convex in $\mathcal W$. Cutting the work requirement by a factor cuts headcount by more than that factor.
\item[(f)] (Scope versus market.) If instead the scope is not fixed and output sells at $P$ per person-year equivalent, the profit-maximizing team size $n^\pi=\tfrac12+(A_g-w/P)/\chi_h$ rises with $A_g$.
\item[(g)] (Agent pods.) With $\chi_a>0$, agents' joint output is at most $(Y_a+\chi_a/2)^2/(2\chi_a)$, and the solo condition also requires $\nu_a\mathcal W_0$ not to exceed it.
\end{enumerate}
\end{proposition}

\noindent\emph{Intuition and example.} Leverage obeys Amdahl's law: the work agents cannot do and the review of the work they do bound how far one person can be multiplied, and an oversight floor binds as hard here as in Corollary~\ref{cor:ext-floor}. Because coordination costs are quadratic, a smaller work requirement removes headcount and coordination load together. With $\mathcal W_0=8$ and $\chi_h=0.03$ the team needs 9.11 people. With $A_g=1.5$, $\bar\nu_a=0.6$ and $\xi^t=0.15$ the requirement falls to 2.85 person-years and the team to 2.94 people. With $A_g=2$, $\nu_a=0.9$ and $\xi^t=0.05$, leverage is 10.5 and the scope is a solo job. Part (f) is the caveat: when the market rather than the work is the bottleneck, AI raises the value of each member and teams grow, the superstar mechanism of \citet{rosen1981economics}. In a field experiment, individuals working with AI matched two-person teams working without it \citep{dellacqua2026cybernetic}.

\begin{assumption}[Knowledge hierarchy]\label{ass:ext-hierarchy}
Problems arrive one per production worker per period, with difficulty $z\sim{\rm Exp}(1)$. A worker who knows $[0,z_w]$ at cost $c_Kz_w$ solves problems with $z\le z_w$ and escalates the rest to a manager, who knows $[0,z_m]$ at cost $c_Kz_m$ and spends $\kappa_c$ of her unit of time on each escalated problem. Wages are $w_w$ and $w_m$, and $z_m$ is set by the quality the firm requires. An AI layer, if present, solves every problem with $z\le z_A$ at negligible cost before humans see it.
\end{assumption}

\begin{proposition}[Spans, layers and the AI layer]\label{prop:ext-hierarchy}
\begin{enumerate}
\item[(a)] The cost per production worker,
\[
C(z_w)=w_w+c_Kz_w+\kappa_ce^{-z_w}(w_m+c_Kz_m),
\]
is minimized at $z_w^*=\ln[\kappa_c(w_m+c_Kz_m)/c_K]$ when $0<z_w^*<z_m$, and the span of control is $\mathrm{sp}^*=1/(\kappa_ce^{-z_w^*})=(w_m+c_Kz_m)/c_K$.
\item[(b)] Cheaper knowledge ($c_K\downarrow$, AI as a source of know-how) raises worker knowledge and widens spans. Cheaper communication ($\kappa_c\downarrow$, AI summaries and shared context) lowers worker knowledge one-for-one in logs. With exponential difficulty it leaves spans unchanged; for a general density $f$ with hazard $\lambda_z$, $d\ln\mathrm{sp}^*/d\ln\kappa_c=-1+\lambda_z/(\lambda_z-\lambda_z'/\lambda_z)$. Cheaper communication therefore widens spans when the hazard is decreasing (heavy-tailed problems) and narrows them when it is increasing.
\item[(c)] (AI layer.) Above an AI layer at $z_A<z_m$, the cost-minimizing human hierarchy is the original one with $z_m$ replaced by $z_m-z_A$. Human headcount per problem falls by the factor $e^{-z_A}$, and spans narrow slightly, to $(w_m+c_K(z_m-z_A))/c_K$.
\item[(d)] (Capability path.) If task lengths are Pareto with tail index $\zeta_P$, so that $z=\zeta_P\ln(d/d_{\min})$, and the AI layer solves tasks with $d\le H(t)$, then the human share of problems is $(H(t)/d_{\min})^{-\zeta_P}$, which falls at the rate $\zeta_P\ln2/T_d$. With $N_w$ production workers the number of layers is about $\ln N_w/\ln\mathrm{sp}^*$, so the AI layer removes about $z_A/\ln\mathrm{sp}^*$ layers.
\end{enumerate}
\end{proposition}

\noindent\emph{Intuition.} AI that makes knowledge cheap decentralizes, and AI that makes communication cheap centralizes, as in \citet{garicano2000hierarchies} and \citet{garicano2006organization}. An AI layer at the bottom leaves a human organization of the same shape, scaled down: by memorylessness, the problems it passes up look like the original problems, shifted. With $\zeta_P=0.5$ and $T_d=207$ days, the human exception load halves every 1.13 years. The tail index is assumed.

\begin{definition}[Residual organization and transformation team]\label{def:ext-org}
With $t_{\rm FTE}$ paid hours per full-time equivalent (2,080, the convention behind annual wages in the BLS Occupational Employment and Wage Statistics) and the optimal modes $x^*_j$:
\begin{enumerate}
\item[(i)] \emph{Residual organization} for occupation $o$: production $\sum_{j\in o}v_jt^{\rm prod}_{x^*_j}/t_{\rm FTE}$, with $t^{\rm prod}=t$ in mode H and $t^U$ in mode U; review $\sum v_j\xi^*_jt^R_j/t_{\rm FTE}$ in mode A (senior), plus approval hours $\sum v_jt_j(\mathsf a^*_j)/t_{\rm FTE}$ under the failure taxonomy of Section~\ref{sec:tasks-failure}; sign-off $\sum v_j\mathrm{lic}_jt^L_j/t_{\rm FTE}$ (licensed); vendor management $\sum v_j\kappa_Ot_j/t_{\rm FTE}$ in mode O; and managers equal to the workers divided by the span $\mathrm{sp}^*$. The senior share is review plus sign-off over the total. Team sizes come from Proposition~\ref{prop:ext-team}.
\item[(ii)] \emph{Transformation team} in wave $k$ of duration $T_k$ (Section~\ref{sec:ext-program}), in FTE:
\begin{itemize}
\item data and integration engineers, $l_{\rm DI}I^{\rm int}_k/(w_{\rm DE}t_{\rm FTE}T_k)$, where $I^{\rm int}_k$ is the wave's integration spend from Section~\ref{sec:tasks} and $l_{\rm DI}$ its labor share;
\item AI and agent engineers, $l_{\rm AI}F_k/(w_{\rm AI}t_{\rm FTE}T_k)+e_{\rm mnt}N^{\mathrm A}_{<k}$, who build the wave and maintain the $N^{\mathrm A}_{<k}$ processes automated earlier;
\item process owners, $e_{\rm PO}n^{\rm re}_k$ for $n^{\rm re}_k$ redesigned processes;
\item change managers, $A^{\rm aff}_k/(T_kN_{\rm CM})$, where $A^{\rm aff}_k$ is the FTE affected and $N_{\rm CM}$ the FTE one change manager supports per year;
\item risk, governance and audit staff, $e_{\rm RG}n^{\rm reg}_k$ for $n^{\rm reg}_k$ regulated AI systems, and domain experts, $e_{\rm SME}n^{\rm build}_k$;
\item staff kept during the parallel run, $\sum_{p\in\mathcal P_k}(\omega_{\rm par}\mathrm{FTE}^0_p+\nu_{\rm rec}\mathrm{FTE}^{\mathrm H}_p)$.
\end{itemize}
\end{enumerate}
\end{definition}

\begin{lemma}[Staffing identities]\label{lem:ext-staffing}
\begin{enumerate}
\item[(a)] Baseline FTE equals residual FTE plus displaced FTE (plus the FTE transferred to vendors in mode O).
\item[(b)] During wave $k$ at ramp point $\tau$, in-house headcount equals baseline headcount, less the savings $\sum(\mathrm{FTE}^0_p-\mathrm{FTE}^1_p)$ of completed waves, plus $\sum_{p\in\mathcal P_k}[\omega_{\rm par}\mathrm{FTE}^0_p+\nu_{\rm rec}\mathrm{FTE}^{\mathrm H}_p-\varrho(\tau)(\mathrm{FTE}^0_p-\mathrm{FTE}^1_p)]$, plus the transformation team. Each wave opens with a headcount hump of $\sum_{p\in\mathcal P_k}(\omega_{\rm par}\mathrm{FTE}^0_p+\nu_{\rm rec}\mathrm{FTE}^{\mathrm H}_p)$ plus the transformation team, which is proportional to the wave's size.
\item[(c)] The pace, in affected FTE a year, is at most the change capacity
\[
K_{\rm chg}:=\min\bigl\{\kappa_{\rm abs}\mathrm{FTE}_{\rm base},\;N_{\rm CM}\mathrm{FTE}^{\rm CM}_{\max}\bigr\},
\]
the smaller of what the organization can absorb and what its change managers can support.
\end{enumerate}
\end{lemma}

In the support example, mode A turns 120.2 FTE of ticket handling into 18.0 FTE of senior review.

\subsection{Program sequencing, parallel running and change capacity}\label{sec:ext-program}

\begin{assumption}[Program]\label{ass:ext-program}
A program migrates $N$ processes. Process $p$ yields the savings flow $W_p$ a year once migrated (the saving $W$ of Assumption~\ref{ass:tasks-transition}, or $W_p(\mathcal J_p)$ of Section~\ref{sec:tasks} for full automation) and affects $a_p$ FTE, and every migration requires a parallel run of length $T_{\rm par}$. Waves $k=1,\dots,n_w$ run in sequence, each for $T_{\rm par}$. Wave $k$ migrates $n_k$ processes and incurs the dual-running cost flow $C^{\rm dr}(n_k)$, with $C^{\rm dr}$ convex, increasing and $C^{\rm dr}(0)=0$. The cost $C^{\rm dr}$ includes the per-process parallel premium $C^{\rm par}$ of Assumption~\ref{ass:tasks-transition} and any congestion; the quadratic case is $C^{\rm dr}(n)=c^{\rm dr}_1n+c^{\rm dr}_2n^2$. Each wave also carries a fixed program overhead $O_w$. Change capacity limits each wave to $\sum_{p\in\mathcal P_k}a_p\le K_{\rm chg}T_{\rm par}$. The program minimizes
\begin{equation}\label{eq:ext-progcost}
\mathcal C^{\rm prog}=\sum_{k=1}^{n_w}\Bigl[T_{\rm par}\,k\sum_{p\in\mathcal P_k}W_p+T_{\rm par}C^{\rm dr}(n_k)+O_w\Bigr],
\end{equation}
the savings forgone while processes wait, dual running and overhead. This is the first-order approximation of the loss in present value when the program is short relative to $1/r$; the engine computes the discounted version.
\end{assumption}

\begin{proposition}[Wave structure]\label{prop:ext-waves}
\begin{enumerate}
\item[(a)] (Order.) Processes enter waves in decreasing order of $W_p$. When processes differ in the capacity $a_p$ they use and capacity binds, the fractional relaxation of the problem without congestion fills earlier waves in decreasing order of $W_p/a_p$.
\item[(b)] (Front-loading.) With identical processes ($W_p=W$), quadratic $C^{\rm dr}$ and $n_w$ waves, the cost-minimizing sizes are
\begin{equation}\label{eq:ext-front}
n_k=\frac N{n_w}+\frac{W}{2c^{\rm dr}_2}\Bigl(\frac{n_w+1}2-k\Bigr),\qquad k=1,\dots,n_w,
\end{equation}
while all are nonnegative: each wave is smaller than the one before by $W/(2c^{\rm dr}_2)$. The cost is then $\mathcal C(n_w)=T_{\rm par}\bigl[WN(n_w+1)/2+c^{\rm dr}_1N+c^{\rm dr}_2N^2/n_w-W^2n_w(n_w^2-1)/(48c^{\rm dr}_2)\bigr]+n_wO_w$.
\item[(c)] (Square-root rule.) With equal waves, the optimal size is $n^*=\sqrt{(WN/2+O_w/T_{\rm par})/c^{\rm dr}_2}$, and $n_w^*=N/n^*$. Few large waves are optimal when savings, program size or overhead are large and congestion is small, and the number of waves grows only with $\sqrt N$.
\item[(d)] (Change capacity.) With identical $a_p=a$, capacity binds when the optimal first wave of (b) exceeds $n_{\rm cap}=K_{\rm chg}T_{\rm par}/a$. The earliest waves are then at capacity and the later ones follow the decrement of (b). The program lasts at least $Na/K_{\rm chg}$, whatever $T_{\rm par}$ is, and capacity has a positive shadow value. Remark~\ref{rem:ext-seccap} combines this cap with decision bandwidth and security review.
\item[(e)] (Big bang.) If $c^{\rm dr}_2=0$, a single wave is optimal whenever capacity allows.
\item[(f)] (Score path.) During wave $k$ the margin changes by $[\sum_{p\text{ in earlier waves}}W_p-C^{\rm dr}(n_k)]/S$, before the one-time costs of Section~\ref{sec:tasks} and the cash cost of capital spent at the start of the wave. The Rule-of-40 dip is deepest in the first wave, at $C^{\rm dr}(n_1)/S$, and the score stays below its baseline while $\sum_{i<k}n_iW<C^{\rm dr}(n_k)$. Few large waves give a deep, short dip, and many small ones a shallow, long one.
\end{enumerate}
\end{proposition}

\noindent\emph{Intuition and example.} Waiting costs the savings of every process not yet migrated, which favors large waves. Running many processes in parallel at once is increasingly costly, which favors small waves. The square-root rule balances the two, as an economic order quantity does. Since early processes save for longer, the first wave should be the largest. With 12 processes saving \$1 million a year each, $T_{\rm par}=0.5$ years, $c^{\rm dr}_1=0.2$, $c^{\rm dr}_2=0.25$ and $O_w=0.5$ (\$ millions), the best integer plan is three waves of 6, 4 and 2, against a square-root rule of 5.3 processes per wave. If each process affects 40 FTE and the firm can absorb 320 affected FTE a year, waves are capped at 4, the plan becomes 4, 4, 4 over 1.5 years, and the cost rises by \$1.0 million; one more process of capacity per wave would save \$0.75 million. On \$500 million of revenue, the unconstrained plan moves the margin by $-2.04$, $+0.24$ and $+1.72$ points in its three waves, and by $+2.4$ points at the end.

\subsection{Risky portfolios, common-mode failure and decision bandwidth}\label{sec:ext-bandwidth}

\begin{assumption}[Portfolio of risky waves]\label{ass:ext-portfolio}
Wave $k$ succeeds with probability $\varpi_k$, aggregated from the process-level probabilities $\varpi_p$ of \eqref{eq:tasks-success}. It is worth $V_k$ on success and loses $C_k$ on failure. As in \eqref{eq:tasks-npv}, a failure costs the up-front outlays and the parallel premium and, if it goes undetected during the parallel run, the stranded costs $C^{\rm str}$ and the rollback loss $L_{\rm roll}$ of Assumption~\ref{ass:tasks-transition} as well. An undetected failure has probability $p_{\rm fail}=(1-\varpi_k)e^{-T_{\rm par}/T_{\rm det}}$, and $\varpi_kV_k-(1-\varpi_k)C_k$ is the expected NPV of Proposition~\ref{prop:tasks-transition}(a). Experience accumulates and decays, with transfer factor $\eta_{\rm tr}\le1$ across business units, and it raises the process-level factor $\varpi^c_p$ of \eqref{eq:tasks-success}, of which team experience is a driver:
\begin{gather*}
E_{k+1}=(1-\delta_E)E_k+\eta_{\rm tr}\times(\text{successes in wave }k),\\
{\rm logit}\,\varpi^c_p(E_k)={\rm logit}\,\varpi^c_p(0)+\upsilon_L\ln(1+E_k)\quad\text{for }p\text{ in wave }k.
\end{gather*}
Waves on a shared platform succeed only if the platform is sound, with probability $p_P$, and independently otherwise. Cost overruns have a power-law upper tail with index $\zeta_T>1$.
\end{assumption}

Knowledge acquired in production depreciates and transfers imperfectly across units \citep{argote1990persistence,darr1995acquisition}, which motivates $\delta_E>0$ and $\eta_{\rm tr}<1$. IT cost overruns follow a power law with tail index about 2.3 above an actual-to-estimate ratio of 2 \citep{flyvbjerg2022empirical}; across 1,471 projects the average overrun was 27\%, and one project in six overran by 200\% on average \citep{flyvbjerg2011why}. These sources do not measure benefit realization; Part I grounds the success probability and the realized share in surveys of transformation outcomes, which are self-reported, so they remain ranges, and the examples below use illustrative values.

\begin{proposition}[Learning, common shocks and tails]\label{prop:ext-portfolio}
\begin{enumerate}
\item[(a)] (Learning and order.) For two processes with discount factor $\delta_d$ per wave and learning gains $\Delta\varpi_j$ (the rise in $j$'s success probability if the other succeeds first),
\begin{equation}\label{eq:ext-order}
EV(1\to2)-EV(2\to1)=(1-\delta_d)(\varpi_1V_1-\varpi_2V_2)+\delta_d(\varpi_1\Delta\varpi_2V_2-\varpi_2\Delta\varpi_1V_1).
\end{equation}
A success shifts the logit of the other process's factor $\varpi^c_j$ by $u:=\upsilon_L\ln(1+\eta_{\rm tr})$, so $\Delta\varpi_j\approx\varpi_j(1-\varpi^c_j)u$, which is $\varpi_j(1-\varpi_j)u$ when the activity and congestion factors of $j$ are close to one. If the two processes share those factors, the learning term has the sign of $\varpi_1-\varpi_2$, so for equal values the likely win goes first. A success $k$ waves earlier shifts the logit by $\upsilon_L\ln[1+\eta_{\rm tr}(1-\delta_E)^{k-1}]$, so learning is worth less across long gaps and across units.
\item[(b)] (Common platform.) For two waves on one platform, each with $\varpi$, $V$ and $C$, and the option to stop after a failure,
\begin{equation}\label{eq:ext-seq}
EV^{\rm seq}-EV^{\rm sim}=\delta_d(1-p_P\varpi)\bigl(-EV_2(p')\bigr)_+-(1-\delta_d)EV_1,\qquad p'=\frac{p_P(1-\varpi)}{1-p_P\varpi},
\end{equation}
where $EV_1=p_P\varpi V-(1-p_P\varpi)C$ and $EV_2(p')=p'\varpi V-(1-p'\varpi)C$. A small first wave on the platform pays when a failure is informative and continuing after one would lose money.
\item[(c)] (Correlation.) With $n_w$ waves on one platform, the number of successes has variance $p_Pn_w\varpi(1-\varpi)+n_w^2\varpi^2p_P(1-p_P)$. The common term grows with $n_w^2$, so waves on one platform do not diversify, and the probability that none succeeds is at least $1-p_P$.
\item[(d)] (Heavy tails.) If wave overruns are independent with a regularly varying tail of index $\zeta_T>1$, then $\Pr(\sum_kX_k>x)\sim n_w\Pr(X>x)$ as $x\to\infty$, so an extreme total overrun comes from a single wave. With $\zeta_T\approx2.3$ the variance is finite but the third moment is not, and normal approximations understate the upper percentiles of program cost. A stage gate that abandons a wave whose overrun exceeds a cap truncates each $X_k$, and with them the tail of the total.
\item[(e)] (Provider concentration.) Let the waves run on $M$ providers, each sound with probability $p_P$ independently, and let provider $m$ host $n^{\rm w}_m\ge1$ of the $n_w$ waves, with concentration $H_n:=\sum_m(n^{\rm w}_m/n_w)^2$. The number of successes has variance $p_Pn_w\varpi(1-\varpi)+n_w^2\varpi^2p_P(1-p_P)H_n$, and the probability that none succeeds is at least $(1-p_P)^M$. The common term is proportional to the concentration $H_n$, and (c) is the case $M=1$. Activity-level failure modes lower each wave's $\varpi$ through Corollary~\ref{cor:tasks-failure-success}; the provider shock is the part they share, and Proposition~\ref{prop:ext-commonmode} treats it once the processes are in operation.
\end{enumerate}
\end{proposition}

\noindent\emph{Example.} With $\varpi=0.80$ and $0.45$, equal values, activity and congestion factors equal to one, $\upsilon_L=1$, $\eta_{\rm tr}=1$ and $\delta_d=0.95$, the learning gains are 0.089 and 0.171, and starting with the likely win is worth 0.109 of one process's value. With a platform sound with probability 0.7, $\varpi=0.9$, $V=1$, $C=0.6$ and $\delta_d=0.95$, a failure lowers the probability that the platform is sound to 0.19, and sequencing raises expected value from 0.816 to 0.911.

\begin{assumption}[Common-mode failure]\label{ass:ext-commonmode}
In operation, the automated processes run on $M$ model providers or model versions. Process $p$ loses $x_p$ a year while its model is unavailable or faulty (savings forgone, manual fallback, errors in flight). The exposure is $X:=\sum_px_p$, provider $m$ carries $X_m$ of it, and $H:=\sum_m(X_m/X)^2$ is the concentration index. Common-mode events at a provider, such as an outage, a regression after a model update or a shared vulnerability, arrive at the rate $\nu^{\rm cm}=\nu^{\rm out}+\nu^{\rm chg}(\kappa_{\rm itgc})$, independently across providers: outages and shared vulnerabilities, which the firm does not control, and regressions introduced by changes to the model, its prompts or its tools, which the firm's change management (pinned versions, testing before promotion, rollback: the IT general controls of Section~\ref{sec:acct-scope}) lowers, so that $\nu^{\rm chg}$ is nonincreasing in the quality $\kappa_{\rm itgc}$ of those controls. The events disrupt every process on it until the event is contained, after the detection lag $T_{\rm mon}$ and the response time $T_{\rm resp}$; write $T_{\rm con}:=T_{\rm mon}+T_{\rm resp}$. An event that disrupts the exposure $y$ for the time $T_{\rm con}$ costs $T_{\rm con}[y+y^2/(2Y_f)]$, where $Y_f$ measures the firm's fallback and incident-response capacity: when many processes fail at once, manual fallback runs short and losses rise more than proportionally. Failures that are independent across processes given the model are treated at the activity level in Section~\ref{sec:tasks-failure}: errors and hallucination under review (Proposition~\ref{prop:tasks-halluc}), automation bias (Proposition~\ref{prop:tasks-bias}), autonomy and misbehavior (Proposition~\ref{prop:tasks-autonomy}) and agent pipelines (Proposition~\ref{prop:tasks-pipeline}). This assumption adds the failures that processes share.
\end{assumption}

Economies of scale and scope in foundation models may push their market toward concentration \citep{korinek2025concentrating}, so many firms will rely on few providers, and a firm that standardizes on one model concentrates its own exposure further. Independence across providers is an assumption. Models from different companies make correlated errors: on a benchmark of multiple-choice questions, two models that are both wrong give the same wrong answer about 60\% of the time, against a third by chance, and more accurate models agree more \citep{kim2025correlated}. That figure is the probability of the same answer given that both are wrong, not the probability that one provider fails when another does, which is not measured; with providers that share base models, cloud regions and training data, the concentration index $H$ is a lower bound on the effective concentration of the exposure. The cost form and its parameters are assumptions; no source we know of measures how the cost of an outage grows with the share of a firm's processes it disrupts.

\begin{proposition}[Common-mode failure, model diversity and governance]\label{prop:ext-commonmode}
Under Assumption~\ref{ass:ext-commonmode}:
\begin{enumerate}
\item[(a)] (Concentration.) The annual cost of common-mode events has mean $\nu^{\rm cm}T_{\rm con}[X+HX^2/(2Y_f)]$ and variance $\nu^{\rm cm}T_{\rm con}^2\sum_m[X_m+X_m^2/(2Y_f)]^2$. The part of the mean that is linear in $X$ does not depend on how the exposure is spread; the rest rises with $H$. Mean and variance both rise when the provider shares become more unequal, in the sense of majorization, and are largest with a single provider, when one event can disrupt the whole exposure.
\item[(b)] (Diversity and redundancy.) Spreading the exposure equally over $M$ providers instead of one saves $\nu^{\rm cm}T_{\rm con}X^2(1-1/M)/(2Y_f)$ a year and divides the variance by at least $M$. Redundancy that fails a disrupted process over to an independent provider, leaving the share $\kappa_{\rm fo}$ of its exposure disrupted, is worth $\nu^{\rm cm}T_{\rm con}[(1-\kappa_{\rm fo})X+(1-\kappa_{\rm fo}^2)HX^2/(2Y_f)]$ a year; unlike diversity, it also cuts the linear part.
\item[(c)] (Scale against diversification.) Let the inference spend with a single provider be $C^{\rm inf}$, with volume discounts of elasticity $b_v\in[0,1)$, so that $M$ equal providers cost $C^{\rm inf}M^{b_v}$, and let each provider add the fixed cost $F_M$ a year for integration, evaluation and revalidation at each model change (the run costs of Section~\ref{sec:tasks}). The annual cost $J(M)=C^{\rm inf}M^{b_v}+MF_M+\nu^{\rm cm}T_{\rm con}[X+X^2/(2Y_fM)]$ is minimized at the smallest $M\ge1$ with
\begin{equation}\label{eq:ext-providers}
\frac{\nu^{\rm cm}T_{\rm con}X^2}{2Y_f}\le M(M+1)\bigl[F_M+C^{\rm inf}\bigl((M+1)^{b_v}-M^{b_v}\bigr)\bigr].
\end{equation}
A single provider is optimal if and only if $\nu^{\rm cm}T_{\rm con}X^2/(4Y_f)\le F_M+C^{\rm inf}(2^{b_v}-1)$. Without volume discounts the optimum is the integer just below or just above $M^\dagger:=\sqrt{\nu^{\rm cm}T_{\rm con}X^2/(2Y_fF_M)}$, a square-root rule. Diversity pays for large exposures, frequent or slow-to-contain events and thin fallback capacity; fixed costs per provider and volume discounts favor a single provider.
\item[(d)] (Staged updates and rollback.) Let a model update be worth $B_{\rm up}$ if it is sound and be a regression with probability $1-p_g$, disrupting what it reaches for $T_{\rm con}$ until it is rolled back. Releasing it first to the share $\eta_c$ of the exposure, and to the rest only if it is sound, delays the rest at a cost of the fraction $\delta_c$ of its value. The expected value $p_gB_{\rm up}[1-(1-\eta_c)\delta_c]-(1-p_g)T_{\rm con}[\eta_cX+(\eta_cX)^2/(2Y_f)]$ is concave in $\eta_c$ and is maximized at
\[
\eta_c^*=\min\Bigl\{1,\ \max\Bigl\{\eta_{\min},\ \frac{Y_f}X\Bigl[\frac{p_gB_{\rm up}\delta_c}{(1-p_g)T_{\rm con}X}-1\Bigr]\Bigr\}\Bigr\},
\]
where $\eta_{\min}$ is the smallest feasible first release. Releasing everything at once is optimal if and only if $p_gB_{\rm up}\delta_c\ge(1-p_g)T_{\rm con}X(1+X/Y_f)$. A staged release is the pilot of Proposition~\ref{prop:tasks-pilot} applied to a model change. Rollback is what limits a regression to $T_{\rm con}$: if a regression instead lasted until a rebuild after $T_{\rm reb}>T_{\rm con}$, keeping a rollback path, such as a pinned earlier model version, would be worth $(1-p_g)(T_{\rm reb}-T_{\rm con})[\eta_cX+(\eta_cX)^2/(2Y_f)]$ per update.
\item[(e)] (Kill switches and decision rights.) A kill switch that cuts the response time from $T_{\rm resp}$ to $T_{\rm ks}$ is worth $\nu^{\rm cm}(T_{\rm resp}-T_{\rm ks})[X+HX^2/(2Y_f)]$ a year. Its value rises with concentration, and the value of diversity rises with the containment time, so fast containment and diversity are substitutes. If business units choose their own exposure $x_u$, through the autonomy they grant agents and the provider they use, and bear event costs in proportion to their exposure, then a unit's marginal cost of exposure on provider $m$ falls short of the firm's by $\nu^{\rm cm}T_{\rm con}(X_m-x_u)/(2Y_f)$. Decentralized choices of exposure are therefore too high. An internal charge of that amount per unit of exposure, evaluated at the firm's optimum, or keeping decision rights over provider choice and autonomy limits at the level that bears the common-mode cost, restores the optimum; with many small units the charge tends to $\nu^{\rm cm}T_{\rm con}X_m/(2Y_f)$. This adds the common-mode loss to the delegation rule of Proposition~\ref{prop:ext-bandwidth}(e) and to the unit-level choice of autonomy in Proposition~\ref{prop:tasks-autonomy}(b), whose expected loss should include the charge.
\item[(f)] (Change management.) The mean of (a) falls by $-\nu^{\rm chg\prime}(\kappa_{\rm itgc})\,T_{\rm con}[X+HX^2/(2Y_f)]$ per unit of control quality $\kappa_{\rm itgc}$, so better change management is worth its marginal cost up to that amount; like a kill switch, it is worth more under concentration and slower containment. Staged release and a pinned rollback path, part (d), are instances. In Proposition~\ref{prop:acct-amin} the same controls lower the common-mode probability of AI control layers, so one investment lowers both operational and misstatement risk.
\end{enumerate}
\end{proposition}

\noindent\emph{Intuition and example.} Independent failures average out across processes; failures that share a model do not, and their cost grows with the share of the firm they disrupt at once because fallback capacity is limited. Diversity attacks the convex part of that cost, redundancy also the linear part, and kill switches and rollback shorten every event. Take \$20 million a year of exposure, two events a year, containment in 0.01 years (under four days) and fallback capacity $Y_f=\$10$ million a year. The expected cost is \$0.80 million a year with one provider and \$0.60 million with two. At \$150{,}000 a year per extra provider, two providers are optimal ($M^\dagger=1.63$), but volume discounts of elasticity 0.1 on \$3 million of inference spend make a single provider optimal. A kill switch that cuts the response from 0.008 to 0.001 years is worth \$0.56 million a year with one provider and \$0.42 million with two. For an update worth \$0.5 million that is a regression one time in ten, with a 10\% cost of delay, the best first release covers 62.5\% of the exposure; if a regression without rollback lasted 0.05 years, the rollback path would be worth \$81{,}000 per update. A business unit that ignores the common-mode cost it imposes on others undervalues exposure by \$0.02 per dollar of exposure a year in this example, a third of the firm's marginal cost of \$0.06. The loss floor $c_{\rm err}$ of \eqref{eq:tasks-cerr} counts expected losses that are independent across activities; at full automation, the expected common-mode cost of part (a), divided by revenue, is a further cost of the same kind in Part I's ceiling, and $c_{\rm err}$ itself must be evaluated at the catch rates of the review that remains, an AI reviewer's net of the errors it shares with the producing model or none, not at human rates (Proposition~\ref{prop:tasks-cerr}(c)).

\begin{assumption}[Decision bandwidth]\label{ass:ext-bandwidth}
The organization has decision bandwidth $\Phi$ per period: leadership attention, governance reviews and approvals. Initiative $i$ uses $e_i$ decision units per period while active, its launch and gates averaged over its duration, so the load is $L=\sum_ie_i$. Congestion lowers success, $\varpi_i=\varpi^0_iq(L/\Phi)$ with $q$ nonincreasing, and raises decision latency $\ell_{\rm dec}(L/\Phi)$. Value erodes at $r_d=r+\delta_V$ while a decision waits, where $\delta_V$ reflects capability, prices and competitors moving on (Sections~\ref{sec:ext-dynamics} and \ref{sec:ext-competition}). Let $Q(\varkappa):=q(\varkappa)e^{-r_d\ell_{\rm dec}(\varkappa)}$.
\end{assumption}

\begin{proposition}[Concurrency, bandwidth and decision rights]\label{prop:ext-bandwidth}
Consider $n$ identical concurrent initiatives, each with load $e$, worth $V$ on success, and with dual-running cost $C^{\rm dr}(n)$.
\begin{enumerate}
\item[(a)] (Load breakpoint.) Expected value is $n\varpi^0VQ(ne/\Phi)-C^{\rm dr}(n)$. If $\varkappa Q(\varkappa)$ is single-peaked, its gross part peaks at the load $\hat \varkappa$ where the elasticity $-\varkappa Q'(\varkappa)/Q(\varkappa)$ equals one. Beyond $\hat \varkappa$, another initiative lowers expected value even without dual-running costs. With convex $C^{\rm dr}$ the optimum lies below $\hat \varkappa\Phi/e$ and satisfies $\varpi^0V[Q(\varkappa)+\varkappa Q'(\varkappa)]=C^{\rm dr\prime}(n)$. With queueing latency $\ell_{\rm dec}=\ell_0/(1-\varkappa)$ and $q\equiv1$, $\hat \varkappa=1+c/2-\sqrt{c+c^2/4}$ with $c=r_d\ell_0$; with quality congestion $q(\varkappa)=e^{-\eta_q(\varkappa-1)_+}$ and no latency, $\hat \varkappa=1/\eta_q$ for $\eta_q<1$, and $\hat \varkappa=1$ otherwise.
\item[(b)] (Pace.) The program migrates at most $\hat \varkappa\Phi/(eT_{\rm par})$ processes a year, so it lasts at least
\[
T_{\rm prog}\;\ge\;\max\Bigl\{\frac{Na}{K_{\rm chg}},\;\frac{NeT_{\rm par}}{\hat \varkappa\Phi}\Bigr\}.
\]
Whichever of change capacity and decision bandwidth is scarcer binds; Remark~\ref{rem:ext-seccap} adds security review.
\item[(c)] (Value of bandwidth.) At the optimum, $d\,EV^*/d\Phi=\bigl[n^*\varpi^0VQ(\varkappa^*)/\Phi\bigr]\,\varepsilon_Q(\varkappa^*)$, the expected portfolio value per decision unit times the elasticity of $Q$. Expanding decision capacity, for example through a transformation office or delegated authority, pays when this exceeds the marginal cost of a decision unit.
\item[(d)] (AI-assisted analysis.) If a decision takes $e=e^{\rm appr}+e^{\rm prep}$, and AI with leverage $A_\Phi$ divides only the preparation part, then bandwidth rises by the factor $(e^{\rm appr}+e^{\rm prep})/(e^{\rm appr}+e^{\rm prep}/(1+A_\Phi))<1+e^{\rm prep}/e^{\rm appr}$. Accountable approval, which stays human, caps the gain, as $a_{\min}$ caps Part I's score.
\item[(e)] (Decision rights.) Let $p_\Phi$ be the shadow price of top-level bandwidth from (c). Delegating initiative $i$ to a lower level pays if and only if $(\varpi^{\rm top}_i-\varpi^{\rm del}_i)V_i<p_\Phi e_i$. Initiatives whose approval must stay with an accountable officer, such as regulated or licensed activities, cannot be delegated and are bandwidth-intensive.
\item[(f)] (Waves.) Replacing $C^{\rm dr}(n)$ by $\tilde C^{\rm dr}(n)=C^{\rm dr}(n)+n\varpi^0V[1-Q(ne/\Phi)]$ gives the wave problem of Proposition~\ref{prop:ext-waves} with congestion. Its results hold on any range of loads on which $\varkappa Q(\varkappa)$ is concave, which makes $\tilde C^{\rm dr}$ convex.
\end{enumerate}
\end{proposition}

\noindent\emph{Intuition and example.} Decision bandwidth behaves like a queue. Near full load, delays and errors grow faster than the number of initiatives, so the best load is well below capacity. With $c=r_d\ell_0=0.125$ (value eroding at 50\% a year while a quarter-year decision waits), the breakpoint is a load of 0.70. With $\varpi^0=0.7$, $V=3$, $e=1$, $\Phi=10$ and $C^{\rm dr}(n)=0.02n^2$, the optimum is 6.8 concurrent initiatives, and one more decision unit is worth 0.78. If approval is 30\% of decision time, AI preparation can at most raise bandwidth 3.3-fold. Delegation follows the logic of \citet{garicano2000hierarchies}: routine decisions go down, exceptions go up, and the price of top-level attention sets the boundary. Attention is a scarce resource of the organization \citep{simon1971designing,ocasio1997towards}, and the architecture of decision rights trades the acceptance of bad projects against the rejection of good ones \citep{sah1986architecture}.

\begin{remark}[Decision bandwidth, security review and change capacity]\label{rem:ext-seccap}
Section~\ref{sec:tasks} models security and compliance review as a second capacity: deployments arrive at the rate $\nu$, are reviewed at capacity $\Phi_{\rm sec}$ a year and wait an exponential time with mean $1/(\Phi_{\rm sec}-\nu)$ (Proposition~\ref{prop:tasks-secreview}). In the setting of Proposition~\ref{prop:ext-bandwidth}, let each of the $n$ concurrent initiatives be active for $T_{\rm par}$ and then deploy, so that $\nu=n/T_{\rm par}$, passing the decision gate and then security review. By Proposition~\ref{prop:tasks-secreview}(d) a deployment keeps the share $Q_{\rm sec}(n):=(\Phi_{\rm sec}-\nu)/(\Phi_{\rm sec}-\nu+r_d)$ of its value, independently of the decision gate, so expected value is $EV(n)=n\varpi^0VQ(ne/\Phi)Q_{\rm sec}(n)-C^{\rm dr}(n)$ for $n<T_{\rm par}\Phi_{\rm sec}$, subject to the change-capacity cap $n\le n_{\rm cap}=K_{\rm chg}T_{\rm par}/a$. Write $\varepsilon_Q$ for the elasticity of $Q$, as in Proposition~\ref{prop:ext-bandwidth}, and $\varepsilon_{\rm sec}(n):=-d\ln Q_{\rm sec}/d\ln n=\nu r_d/[(\Phi_{\rm sec}-\nu)(\Phi_{\rm sec}-\nu+r_d)]$, which increases with $n$.
\begin{enumerate}
\item[(a)] (Binding constraint.) If $\varepsilon_Q(ne/\Phi)+\varepsilon_{\rm sec}(n)$ is continuous and strictly increasing in $n$, the gross value $n\varpi^0VQ\,Q_{\rm sec}$ peaks at the unique $\hat n_2$ at which the two elasticities sum to one. It lies below both single-capacity breakpoints, $\hat n_2\le\min\{\hat\varkappa\Phi/e,\,T_{\rm par}\hat\nu_{\rm sec}\}$, strictly below the first and, when $\varepsilon_Q>0$ there, below the second, where $\hat\nu_{\rm sec}:=\Phi_{\rm sec}+r_d-\sqrt{r_d^2+r_d\Phi_{\rm sec}}<\Phi_{\rm sec}$ solves $\varepsilon_{\rm sec}=1$. With congestion in both reviews, neither capacity alone sets the limit: each lowers the breakpoint set by the other. The pace bound of Proposition~\ref{prop:ext-bandwidth}(b) becomes $T_{\rm prog}\ge\max\{Na/K_{\rm chg},\,NT_{\rm par}/\hat n_2\}$, and deployments beyond $\hat\nu_{\rm sec}$ a year lower the value of the deployment flow, a tighter limit than the capacity $\Phi_{\rm sec}$ of Proposition~\ref{prop:tasks-secreview}(d).
\item[(b)] (Shadow values.) At the optimum $n^*$, with multiplier $\zeta\ge0$ on the change-capacity cap, $\partial EV^*/\partial\Phi=[n^*\varpi^0VQ\,Q_{\rm sec}/\Phi]\,\varepsilon_Q(\varkappa^*)$, which is Proposition~\ref{prop:ext-bandwidth}(c) with $V$ replaced by $VQ_{\rm sec}(n^*)$. The value of a unit of review capacity per year, $T_{\rm par}^{-1}\partial EV^*/\partial\Phi_{\rm sec}$, is \eqref{eq:tasks-secvalue} at $\nu=n^*/T_{\rm par}$ and $V_{\rm dep}=\varpi^0VQ(\varkappa^*)$, the value of a deployment net of decision congestion; given $n^*$, the review capacity worth buying follows the square-root rule of Proposition~\ref{prop:tasks-secreview}(c) at these values. A unit of change capacity, one affected FTE a year, is worth $\zeta T_{\rm par}/a$, where $\zeta=\varpi^0VQ\,Q_{\rm sec}(1-\varepsilon_Q-\varepsilon_{\rm sec})-C^{\rm dr\prime}(n_{\rm cap})>0$ if the cap binds and $\zeta=0$ otherwise. With smooth congestion both review capacities have positive shadow values whether or not change capacity binds, and the firm should expand the capacity whose shadow value most exceeds its marginal cost.
\item[(c)] (Hard caps.) If instead each review runs without congestion up to its capacity and cannot exceed it, concurrency is limited by $\bar n:=\min\{K_{\rm chg}T_{\rm par}/a,\,\Phi/e,\,T_{\rm par}\Phi_{\rm sec}\}$, and Proposition~\ref{prop:ext-waves}(d) holds with $n_{\rm cap}=\bar n$. In its fractional relaxation the program cost $\mathcal C^*(\bar n)$ is convex and nonincreasing, with $-\mathcal C^{*\prime}(\bar n)=\sum_k\zeta_k$, the sum of the multipliers on the wave caps. Only the capacity that attains the minimum has value: $T_{\rm par}\sum_k\zeta_k/a$ per affected FTE a year, $\sum_k\zeta_k/e$ per decision unit, or $T_{\rm par}\sum_k\zeta_k$ per review a year. The other two are worth nothing at the margin, and at a tie raising one capacity alone is worth nothing. Hard capacities are perfect complements: expanding the scarcest pays only until the next one binds.
\end{enumerate}
\end{remark}

\noindent In the bandwidth example above, with $T_{\rm par}=0.5$ years and capacity for 20 security reviews a year, the breakpoint is 7.0 concurrent initiatives for decision bandwidth alone and 8.6 for security review alone, but 6.8 for the two together. The optimum falls from 6.8 to 6.5 initiatives, a decision unit is worth 0.61 instead of 0.78, and a unit of review capacity is worth 0.17 a year. When capacities are hard caps, the scarcest one alone has value: in the program example of Proposition~\ref{prop:ext-waves}, with four waves and change capacity binding at four processes a wave, an extra affected FTE a year of capacity saves \$12,500 in the relaxed problem, while extra decision bandwidth or review capacity saves nothing; if review capacity is the scarcest, an extra review a year saves \$500,000.

\subsection{Regime shocks: leadership, vision, reorganization and context}\label{sec:ext-shocks}

Transformations last years, and the people, priorities and structures that start them change before they finish.

\begin{assumption}[Regime shocks]\label{ass:ext-shocks}
Four shocks arrive as Poisson processes. Sponsor or leadership changes arrive at rate $\nu^{\rm L}$, so expected tenure is $1/\nu^{\rm L}$. Vision overhauls arrive at rate $\nu^{\rm V}$, which may jump after a leadership change. Reorganizations arrive at rate $\nu^{\rm R}$, and departures of each of the $n_K$ people who hold critical process knowledge at rate $\nu^{\rm K}$. A process is \emph{in flight} from its start until it is embedded at $T_{\rm emb}$, its payback date including the parallel run (Proposition~\ref{prop:tasks-transition}(c)), after which its savings are durable.
\begin{enumerate}
\item[(i)] A leadership change keeps the program with probability $\kappa_{\rm keep}$. Otherwise the program is rescoped or cancelled, which, like a vision overhaul, disrupts every process in flight. The disruption rate during flight is $\nu_D:=\nu^{\rm L}(1-\kappa_{\rm keep})+\nu^{\rm V}$, where $\nu^{\rm V}$ is the average rate including overhauls triggered by kept leaders. A disrupted process reverts to its legacy mode. Its built assets carry the share $s^{\rm reuse}$ of their cost $F+I$ into the next plan, and re-planning costs $C^{\rm rep}$, so the value recovered is $V^{\rm rec}:=s^{\rm reuse}(F+I)-C^{\rm rep}$. A new leader also resets the sponsorship driver of $\varpi_p$ in \eqref{eq:tasks-success}.
\item[(ii)] A reorganization freezes every initiative in flight for $T_{\rm frz}$, during which decision bandwidth and change capacity are unavailable. A frozen wave must be re-mobilized at cost $O^{\rm re}_w$, and handoff costs and process owners may change.
\item[(iii)] The organization's context stock $\mathfrak m_t\in[0,1]$ recovers at the relearning rate $\eta_K$. Each departure lowers it by $(1-\kappa_{\rm doc})/n_K$, where $\kappa_{\rm doc}$ is the share of a person's context captured by documentation, context engines and data platforms. Lost context costs $L^{\rm ctx}$ a year per unit, through higher review shares, more errors and lower readiness. In the additive approximation, $d\mathfrak m=\eta_K(1-\mathfrak m)\,dt-[(1-\kappa_{\rm doc})/n_K]\,dN^{\rm K}_t$, where $N^{\rm K}$ counts departures, at rate $n_K\nu^{\rm K}$.
\end{enumerate}
\end{assumption}

For chief executives of large U.S.\ firms the turnover hazard is about 0.16 a year \citep{kaplan2012how}, and forced turnover rises after industry downturns \citep{jenter2015ceo}, so leadership changes can coincide with other shocks. Successors make fewer strategic changes when the former chief executive stays on as board chair \citep{quigley2012when}, which is consistent with $\nu^{\rm V}$ jumping after a leadership change and is one governance arrangement that acts as a commitment device. Tenure data for chief information, data or AI officers are practitioner data. No evidence we know of links reorganizations to the timing of initiatives in flight, so $T_{\rm frz}$ and $O^{\rm re}_w$ are assumptions. We treat organizational memory, a concept developed by \citet{walsh1991organizational}, as a stock held partly by people and partly in documentation and systems; $\kappa_{\rm doc}$ is the share of a person's context held outside the person.

Write $r_D:=r+\nu_D$ for the effective discount rate of work in flight. Let $V^{\rm pay}(\tau)$ be the time-0 value of a process's payoff if it is disrupted at $\tau<T_{\rm emb}$: savings $W$ a year from $T_{\rm par}$ until $\tau$, then $V^{\rm rec}$. Let $V^{\rm pay}(\infty)=We^{-rT_{\rm par}}/r$ be its value if it is never disrupted.

\begin{proposition}[Program value under regime shocks]\label{prop:ext-shocks}
Consider a process that costs $F+I$ at date 0, saves $W$ a year from $T_{\rm par}$ (the risk-adjusted saving $\varpi_p\bar b_pW$ of \eqref{eq:tasks-npv} when execution risk matters), and is in flight until $T_{\rm emb}\ge T_{\rm par}$.
\begin{enumerate}
\item[(a)] (Value.) Its expected value is
\begin{equation}\label{eq:ext-Vshock}
V(\nu_D)=-(F+I)+W\Bigl[\frac{e^{-r_DT_{\rm par}}-e^{-r_DT_{\rm emb}}}{r_D}+\frac{e^{-r_DT_{\rm emb}}}{r}\Bigr]+V^{\rm rec}\,\frac{\nu_D}{r_D}\bigl(1-e^{-r_DT_{\rm emb}}\bigr).
\end{equation}
It is decreasing in $\nu_D$, and hence in $\nu^{\rm L}$ and $\nu^{\rm V}$, and increasing in $\kappa_{\rm keep}$, if $rV^{\rm rec}\le W$ and either $V^{\rm rec}\le0$ or $T_{\rm par}=0$; strictly so if $rV^{\rm rec}<W$. To first order, $V(\nu_D)-V(0)\approx-\nu_D\int_0^{T_{\rm emb}}[V^{\rm pay}(\infty)-V^{\rm pay}(\tau)]\,d\tau$: the loss is the hazard times the area of value at risk, which grows with the payback window. Payback-heavy programs therefore bear most of the loss.
\item[(b)] (Tenure breakpoint.) The probability of embedding before a disruption is $e^{-\nu_DT_{\rm emb}}$. If the value secured at embedding grows in proportion to the window, $V_{\rm emb}=\bar v\,T_{\rm emb}$, the expected discounted value $\bar vT_{\rm emb}e^{-r_DT_{\rm emb}}$ peaks at $T_{\rm emb}=1/r_D$ and falls for longer windows. The breakpoint is therefore $r_DT_{\rm emb}=1$, or, with $r$ and $\nu^{\rm V}$ negligible, expected tenure $1/\nu^{\rm L}=(1-\kappa_{\rm keep})T_{\rm emb}$. When tenure is shorter, programs should be split into wins that embed within $1/r_D$.
\item[(c)] (Sequencing.) If processes run one after another and a disruption cancels the process in flight and all later ones, it is optimal to run them in decreasing order of $V_{{\rm emb},i}/(e^{r_D\tau_i}-1)$, where $\tau_i$ is process $i$'s window. For a quick win ($V_q,\tau_q$) and a long process ($V_l,\tau_l$) with $\tau_q<\tau_l$ and $V_l/\tau_l>V_q/\tau_q$, there is a unique rate $r_D^*$ above which the quick win goes first.
\item[(d)] (Vision-agnostic investment.) $\partial V/\partial s^{\rm reuse}=(F+I)(\nu_D/r_D)(1-e^{-r_DT_{\rm emb}})$, which is positive and increasing in $\nu_D$. Between two otherwise identical investments, the more modular one gains $(s^{\rm reuse}_1-s^{\rm reuse}_2)(F+I)(\nu_D/r_D)(1-e^{-r_DT_{\rm emb}})$, and the gain rises with the hazards.
\end{enumerate}
\end{proposition}

\noindent\emph{Intuition and example.} A disruption destroys value only while a process is in flight, so the hazards act like a higher discount rate on work not yet embedded. Take a process that costs \$10 million, saves \$4 million a year after a six-month parallel run and embeds at three years, with $r=0.10$ and $V^{\rm rec}=0$. It is worth \$28.0 million without shocks and \$7.2 million at $\nu_D=0.3$. With $\kappa_{\rm keep}=0.5$ and $\nu^{\rm V}=0.05$, a three-year window is past the breakpoint once sponsor tenure falls below 2.7 years. A quick win worth \$2 million that embeds in nine months goes before a \$10 million process that takes three years as soon as $r_D$ exceeds 0.19. Part (d) links to the big push of Proposition~\ref{prop:tasks-readiness}(d). Data integration and platforms are reused by any later plan ($s^{\rm reuse}$ near one), so their value is almost immune to vision overhauls, while vision-specific builds are not. Under high hazards the readiness-first sequence of the core section is therefore also the robust one, provided each vision-specific build that follows embeds within $1/r_D$.

\begin{proposition}[Context continuity, commitment and interactions]\label{prop:ext-continuity}
\begin{enumerate}
\item[(a)] (Context continuity.) In the stationary state of Assumption~\ref{ass:ext-shocks}(iii),
\[
\mathbb E[\mathfrak m]=1-\frac{\nu^{\rm K}(1-\kappa_{\rm doc})}{\eta_K},\qquad \mathrm{Var}(\mathfrak m)=\frac{\nu^{\rm K}(1-\kappa_{\rm doc})^2}{2\eta_Kn_K}.
\] Raising $\kappa_{\rm doc}$ by $\Delta\kappa$ is worth $L^{\rm ctx}\nu^{\rm K}\Delta\kappa/(\eta_Kr)$: the insurance value of documentation and context engines. Concentration of knowledge in few people (small $n_K$) leaves the mean unchanged but raises the variance, and documentation lowers the variance quadratically.
\item[(b)] (Commitment devices.) Suppose an overhaul during a window of length $\tau$, which happens with probability $P_V=1-e^{-\nu^{\rm V}\tau}$, is informative with probability $p_{\rm inf}$. An informative overhaul avoids a continuation loss $V^->0$; an arbitrary one destroys a continuation value $V^+>0$. Either way the firm recovers $V^{\rm rec}$. A device that rules out overhauls (a charter, board-level targets, multi-year funding) is worth $P_V[(1-p_{\rm inf})(V^+-V^{\rm rec})-p_{\rm inf}(V^-+V^{\rm rec})]$, the arbitrary-overhaul risk it removes less the abandonment option it gives up. A gated device that allows abandonment only on verifiable evidence, as in the stage gate of Proposition~\ref{prop:tasks-pilot}, is worth $P_V(1-p_{\rm inf})(V^+-V^{\rm rec})$, at least as much. The same holds for devices that raise $\kappa_{\rm keep}$, with $\nu^{\rm L}(1-\kappa_{\rm keep})$ in place of $\nu^{\rm V}$.
\item[(c)] (Bandwidth and change capacity.) Freezes make bandwidth and change capacity available a fraction $1/(1+\nu^{\rm R}T_{\rm frz})$ of the time, so Proposition~\ref{prop:ext-bandwidth}(b) holds with $\Phi$ and $K_{\rm chg}$ divided by $1+\nu^{\rm R}T_{\rm frz}$.
\item[(d)] (Waves.) To first order in $\nu^{\rm R}$, freezes add $\nu^{\rm R}T_{\rm frz}(\mathcal C^{\rm prog}-n_wO_w)+\nu^{\rm R}T_{\rm par}n_wO^{\rm re}_w$ to the program cost \eqref{eq:ext-progcost}. With equal waves the optimal size becomes
\[
n^*=\sqrt{\Bigl[\frac{WN}2+\frac{O_w+\nu^{\rm R}T_{\rm par}O^{\rm re}_w}{(1+\nu^{\rm R}T_{\rm frz})T_{\rm par}}\Bigr]\Big/c^{\rm dr}_2}.
\] Freezes therefore favor fewer, larger waves if and only if $T_{\rm par}O^{\rm re}_w>T_{\rm frz}O_w$, that is, when re-mobilizing a frozen wave is costly relative to the normal overhead per wave. If reorganizations hit one business unit at a time and the program runs as $n_s$ decoupled streams, a freeze stops only its own unit's stream. The expected cost of freezes is then $1/n_s$ of what it is when every freeze stops the whole program.
\end{enumerate}
\end{proposition}

\noindent\emph{Intuition and example.} Context is a stock that departures deplete and relearning rebuilds. Documentation and context engines lower both the expected loss and its variance, and the same platforms raise $\kappa_{\rm ctx}$ in Assumption~\ref{ass:ext-capability}, a double dividend. With $\nu^{\rm K}=0.15$, $\eta_K=1$ and $L^{\rm ctx}=\$5$ million a year, raising $\kappa_{\rm doc}$ from 0.2 to 0.7 lifts the mean context stock from 0.88 to 0.955 and is worth \$3.75 million. With three knowledge holders the standard deviation of the stock falls from 0.13 to 0.05. A commitment device is an option written against the firm's own future leaders. It is worth having when overhauls are mostly arbitrary, and it is best gated: a binding commitment also invites escalation of commitment to a failing course \citep{staw1976knee,staw1981escalation}, and investors value the abandonment option it removes \citep{berger1996investor,trigeorgis1996real}. With $P_V=0.39$ (a two-year window at $\nu^{\rm V}=0.25$), $V^+=10$, $V^-=4$, $V^{\rm rec}=1$ and $p_{\rm inf}=0.3$, a binding charter is worth 1.89 and a gated one 2.48. The claim that reorganization freezes favor fewer waves holds only under the condition in (d): a freeze also delays and prolongs whatever is in flight, which by itself favors smaller waves. Decoupling helps in either case, since a freeze in one unit then stops only that unit's stream. That knowledge acquired in operation depreciates, and transfers only partly across units, is documented by \citet{argote1990persistence} and \citet{darr1995acquisition}.

\subsection{Breakpoints}\label{sec:ext-breaks}

\begin{proposition}[Breakpoints in this section]\label{prop:ext-breaks}
The engine reports the distance of each activity, process or program to the following thresholds.
\begin{enumerate}
\item[(a)] (Organizational.) $n^*=1$ if and only if $\mathcal W_0\le\Lambda$ (Proposition~\ref{prop:ext-team}(c)). The solo date solves $\mathcal W(t)=1$, with $\bar\nu_a(t)$ from Proposition~\ref{prop:ext-feasible} and the paths of $A_g$ and $\xi^t$; it is unique if $\mathcal W$ is decreasing.
\item[(b)] (Market.) $P_I>\hat c_E(t)+p_{\rm sw}$, first at $T_M$ (Proposition~\ref{prop:ext-should}(b)).
\item[(c)] (Regulatory class.) Moving activity $j$ from class $c$ to $c'$ raises $\tilde c_{\mathrm A,j}$ by the rise in its quality cost when the floor $\xi^{\min}(c')$ exceeds the current share $\xi^*_j$, plus any conformity cost per unit $c^K_{\rm conf}$:
\[
J_j=\bigl[c^Q_j(\xi^{\min}(c'))-c^Q_j(\xi^*_j)\bigr]_++(1+\tau_A)\Upsilon\,c^K_{\rm conf}\;\le\;\hat w_jt^R_j\bigl(\xi^{\min}(c')-\xi^*_j\bigr)_++(1+\tau_A)\Upsilon\,c^K_{\rm conf},
\]
with $c^Q_j$ evaluated at $\hat w_j$; the bound holds because the forced review also catches errors. The activity leaves the universe if $0<\mathcal B_j\le J_j$ and no other AI mode wins; class 3 removes it. Where the class depends on the AI's role, as when Article 6(3) of the EU AI Act exempts some assistive or preparatory uses of Annex III systems \citep{eu2024regulation}, the boundary falls between modes U and A: the best mode switches from A to U, and value falls by $\min\{J_j,\tilde c_{\mathrm U,j}-\tilde c_{\mathrm A,j}\}$.
\item[(d)] (Oversight floor.) The floor at which $j$ leaves the universe solves $c^Q_j(\xi^{\min*})-c^Q_j(\xi^*_j)=\mathcal B_j$, so $\xi^{\min*}\ge\xi^*_j+\mathcal B_j/(\hat w_jt^R_j)$, with equality when the expected error loss does not depend on the share reviewed.
\item[(e)] (Timing.) Feasibility $T^F_j$ (Proposition~\ref{prop:ext-feasible}), adoption $T^*_p$ (Proposition~\ref{prop:ext-timing}), and contract expiries as focal dates.
\item[(f)] (Program.) Change capacity binds when $n_{\rm cap}<n_1^*$ (Proposition~\ref{prop:ext-waves}(d)), decision load passes its breakpoint at $\hat\varkappa$ (Proposition~\ref{prop:ext-bandwidth}(a)), and with security review concurrency passes its breakpoint at $\hat n_2$, while hard caps bind at $\bar n$ (Remark~\ref{rem:ext-seccap}).
\item[(g)] (Tenure.) $r_DT_{\rm emb}=1$: with $r$ and $\nu^{\rm V}$ negligible, expected sponsor tenure equal to $(1-\kappa_{\rm keep})$ times payback plus parallel run (Proposition~\ref{prop:ext-shocks}(b)); and the sequencing switch $r_D^*$ (Proposition~\ref{prop:ext-shocks}(c)).
\item[(h)] (Freezes.) $T_{\rm par}O^{\rm re}_w=T_{\rm frz}O_w$, the boundary between freezes favoring fewer or more waves (Proposition~\ref{prop:ext-continuity}(d)).
\item[(i)] (Providers.) A second provider pays when $\nu^{\rm cm}T_{\rm con}X^2/(4Y_f)>F_M+C^{\rm inf}(2^{b_v}-1)$, and further ones by \eqref{eq:ext-providers}; a model update should be staged when $p_gB_{\rm up}\delta_c<(1-p_g)T_{\rm con}X(1+X/Y_f)$ (Proposition~\ref{prop:ext-commonmode}(c) and (d)).
\end{enumerate}
\end{proposition}

\begin{table}[t]
\centering
\caption{What limits the universe available to automate, and what moves each limit}\label{tab:ext-limits}
\footnotesize
\setlength{\tabcolsep}{4pt}
\begin{tabular}{@{}p{3.2cm}p{4.9cm}p{4.2cm}p{2.5cm}@{}}
\toprule
Gap (Lemma~\ref{lem:ext-nested}) & Driven by & Closes over time through & Unaffected by\\
\midrule
Capability, $1-\Sigma^{\rm fea}$ & horizon, reliability, context, physical lag, readiness & $T_d$, context engines, world models & prices, energy, $\Theta$\\
Rules, $\Sigma^{\rm fea}-\Sigma^{\rm per}$ & class 3, licensing, exclusions & changes in law or commitments only & capability, prices\\
Preferences, $\Sigma^{\rm per}-\Sigma^{\rm pre}$ & floors, sign-off, $\omega$, $\vartheta$, $\delta^G$, $\delta^\lambda$, errors & reliability (errors, review), $\Theta$ & AI prices, energy\\
Cost, $\Sigma^{\rm pre}-\Sigma^{\rm eco}$ & inference price, build cost, $\tau_A$, $\Upsilon$ & $\delta_\pi$, $\delta_F$, energy (small) & preferences\\
\bottomrule
\end{tabular}
\end{table}
